\appendix

\section{Proof Details}

\subsection{Unit-box coefficient envelopes and chain derivatives}
\label{app:coefficient-envelope}

\begin{lemma}[Literal unit-box coefficient envelope]
\label{lem:unit-box-envelope}
Let $m\geq1$ and $d\geq0$ be integers, and let
\[
S(z)=\sum_{\nu\in\mathcal A}a_\nu z^\nu,
\qquad
\mathcal A\subseteq
\{\nu\in\mathbb N_0^m:|\nu|_1\leq d\}
\]
have finite coefficient support.  Then, for every $z\in[-1,1]^m$,
\[
|S(z)|\leq\|\operatorname{coeff}(S)\|_1,
\qquad
|\partial_{z_r}S(z)|
\leq d\,\|\operatorname{coeff}(S)\|_1
\quad(1\leq r\leq m).
\]
These statements include $d=0$, constant polynomials, and the zero
polynomial without assigning a degree to the zero polynomial.
\end{lemma}

\begin{proof}
For $z\in[-1,1]^m$, every monomial satisfies $|z^\nu|\leq1$, and hence
\[
|S(z)|
\leq\sum_{\nu\in\mathcal A}|a_\nu||z^\nu|
\leq\sum_{\nu\in\mathcal A}|a_\nu|
=\|\operatorname{coeff}(S)\|_1.
\]
For a fixed coordinate $r$, termwise differentiation gives
\[
\partial_{z_r}S(z)
=\sum_{\substack{\nu\in\mathcal A\\\nu_r\geq1}}
 \nu_r a_\nu z^{\nu-e_r}.
\]
Since $\nu_r\leq|\nu|_1\leq d$ and
$|z^{\nu-e_r}|\leq1$,
\[
\begin{aligned}
|\partial_{z_r}S(z)|
&\leq
\sum_{\substack{\nu\in\mathcal A\\\nu_r\geq1}}
 \nu_r|a_\nu||z^{\nu-e_r}|\\
&\leq d
\sum_{\substack{\nu\in\mathcal A\\\nu_r\geq1}}|a_\nu|
\leq d\,\|\operatorname{coeff}(S)\|_1.
\end{aligned}
\]
No monomial-count factor occurs.  If $d=0$, the differentiated sum is empty;
if $S=0$, both inequalities are $0\leq0$.
\end{proof}

\begin{proposition}[Coordinate derivative envelope]
\label{prop:coordinate-envelope}
Under Assumptions~\ref{assump:balcan-common-chain} and
\ref{assump:anchored-unit-range}, define
\[
D_*:=\Delta B_Q(1+qB_P).
\]
Then $G_i\in C^1([-1,1])$ and
\[
|G_i'(x)|\leq D_*
\qquad
(x\in[-1,1],\ 1\leq i\leq N).
\]
For $q=0$, this is $|G_i'|\leq\Delta B_Q$ with $M=B_P=0$.
For fixed $B_P$, the bound has degree-zero dependence on $M$.
\end{proposition}

\begin{proof}
Fix $x\in[-1,1]$ and $i\leq N$, and write
\[
z(x):=(x,\eta_1(x),\ldots,\eta_q(x)).
\]
When $P_j$ is evaluated at $z(x)$ it ignores later chain coordinates, as in
Assumption~\ref{assump:balcan-common-chain}.  The unit-range condition gives
$z(x)\in[-1,1]^{q+1}$; for $q=0$, this says $z(x)=(x)\in[-1,1]$.

Because $\deg Q_i\leq\Delta$, Lemma~\ref{lem:unit-box-envelope} gives
\[
|\partial_xQ_i(z(x))|\leq\Delta B_Q,
\qquad
|\partial_{y_j}Q_i(z(x))|\leq\Delta B_Q
\quad(1\leq j\leq q).
\]
For $q\geq1$, the value part of that lemma and the definition of $B_P$ give
\[
|P_j(z(x))|
\leq\|\operatorname{coeff}(P_j)\|_1
\leq B_P.
\]
The constant is one because every monomial has magnitude at most one on the
actual unit box; neither $M$ nor a monomial count is introduced.

The multivariable chain rule \citep{rudin1976principles} and the chain
equations yield
\[
\begin{aligned}
G_i'(x)
&=\partial_xQ_i(z(x))
 +\sum_{j=1}^q\partial_{y_j}Q_i(z(x))\eta_j'(x)\\
&=\partial_xQ_i(z(x))
 +\sum_{j=1}^q\partial_{y_j}Q_i(z(x))P_j(z(x)).
\end{aligned}
\]
Consequently,
\[
\begin{aligned}
|G_i'(x)|
&\leq |\partial_xQ_i(z(x))|
 +\sum_{j=1}^q
 |\partial_{y_j}Q_i(z(x))|\,|P_j(z(x))|\\
&\leq\Delta B_Q+\sum_{j=1}^q\Delta B_QB_P\\
&=\Delta B_Q+q\Delta B_QB_P
=\Delta B_Q(1+qB_P)=D_*.
\end{aligned}
\]
The calculation holds on the interval interior.  Every term is continuous on
$[-1,1]$, so the same identity and estimate hold at the endpoints by the
corresponding one-sided limits.

When $q=0$, both sums are empty and
\[
G_i'(x)=Q_i'(x),
\qquad
|G_i'(x)|\leq\Delta B_Q.
\]
If $\Delta=0$, every output is constant and both sides vanish.  A constant or
zero $Q_i$ has zero derivative even when the global $\Delta$ is positive.  If
$M=0$ with $q\geq1$, each nonzero $P_j$ is constant and the same value bound
applies; a zero $P_j$, or $B_P=0$, removes its chain contribution.  These
cases introduce no additional convention or constant.
\end{proof}

\subsection{Euclidean normalization and projective speed}
\label{app:projective-speed}

\begin{lemma}[Derivative of a normalized Euclidean curve]
\label{lem:normalized-curve}
Let $J=[a,b]$ be nondegenerate and let
$v\in C^1(J;\mathbb R^N)$ satisfy $v(t)\neq0$ on $J$.  Define
\[
r(t):=\|v(t)\|_2,
\qquad
\gamma(t):=\frac{v(t)}{r(t)}.
\]
Then $r,\gamma\in C^1(J)$ under the one-sided endpoint convention, and
\[
r'(t)=\gamma(t)^{\mathsf T}v'(t),
\qquad
\gamma'(t)
=\frac{(I_N-\gamma(t)\gamma(t)^{\mathsf T})v'(t)}{\|v(t)\|_2}.
\]
The matrix $I_N-\gamma\gamma^{\mathsf T}$ is the Euclidean orthogonal
projector onto $\gamma^\perp$, so
\[
\|I_N-\gamma(t)\gamma(t)^{\mathsf T}\|_{2\to2}\leq1,
\qquad
\|\gamma'(t)\|_2\leq\frac{\|v'(t)\|_2}{\|v(t)\|_2}.
\]
For $N=1$, the projector is zero and $\gamma$ is constant.
\end{lemma}

\begin{proof}
Nonvanishing gives $r>0$.  Differentiating
$r(t)^2=v(t)^{\mathsf T}v(t)$ gives
\[
2r(t)r'(t)=2v(t)^{\mathsf T}v'(t),
\qquad
r'(t)=\frac{v(t)^{\mathsf T}v'(t)}{r(t)}
=\gamma(t)^{\mathsf T}v'(t).
\]
Differentiating $v/r$ therefore yields
\[
\begin{aligned}
\gamma'(t)
&=\frac{v'(t)}{r(t)}-\frac{v(t)r'(t)}{r(t)^2}\\
&=\frac{v'(t)-\gamma(t)\gamma(t)^{\mathsf T}v'(t)}{r(t)}\\
&=\frac{(I_N-\gamma(t)\gamma(t)^{\mathsf T})v'(t)}{\|v(t)\|_2}.
\end{aligned}
\]
The right sides are continuous, so the interior identities extend to the
endpoints by one-sided limits.

Put $P=I_N-\gamma\gamma^{\mathsf T}$.  Since $\|\gamma\|_2=1$,
\[
P^{\mathsf T}=P,
\qquad
P^2=I_N-2\gamma\gamma^{\mathsf T}
 +\gamma(\gamma^{\mathsf T}\gamma)\gamma^{\mathsf T}=P.
\]
Also $P\gamma=0$, while $Pz=z$ for $z\perp\gamma$, so $P$ is the stated
orthogonal projector.  For every $z\in\mathbb R^N$,
\[
\|Pz\|_2^2=z^{\mathsf T}P^2z
=\|z\|_2^2-(\gamma^{\mathsf T}z)^2
\leq\|z\|_2^2.
\]
This proves the contraction and the derivative bound.  If $N=1$, continuity
and nonvanishing prevent the unit scalar $\gamma$ from changing sign, and
algebraically $P=1-\gamma^2=0$.
\end{proof}

\begin{proposition}[Anchored projective-speed certificate]
\label{prop:projective-speed}
Under Assumptions~\ref{assump:parameter-regime} and
\ref{assump:anchored-unit-range}, and
Proposition~\ref{prop:coordinate-envelope}, define
\[
D_*:=\Delta B_Q(1+qB_P),
\qquad
\gamma_G(x):=\frac{G(x)}{\|G(x)\|_2}.
\]
Then for every $x\in[-1,1]$ and $\theta\in\Theta$,
\[
G_1(x)=F_1(\theta)=1,
\qquad
\|G(x)\|_2,\|F(\theta)\|_2\geq1.
\]
Both normalized curves are globally $C^1$, and
\[
\gamma_G'(x)
=\frac{(I_N-\gamma_G(x)\gamma_G(x)^{\mathsf T})G'(x)}{\|G(x)\|_2},
\]
\[
\gamma_F(\theta)=\gamma_G\!\left(\frac{\theta-c}{h}\right),
\qquad
\gamma_F'(\theta)=\frac1h\gamma_G'\!\left(\frac{\theta-c}{h}\right).
\]
In particular,
\[
\Gamma_{\mathrm{proj}}(F)
\leq\frac{\sqrt N\,D_*}{h}
=\frac{\sqrt N\,\Delta B_Q(1+qB_P)}{h}.
\]
\end{proposition}

\begin{proof}
The literal anchor gives
\[
G_1(x)=Q_1(x,\eta_1(x),\ldots,\eta_q(x))=1,
\qquad
F_1(\theta)=G_1(x(\theta))=1.
\]
Thus
\[
\|G(x)\|_2^2=1+\sum_{i=2}^N|G_i(x)|^2\geq1,
\qquad
\|F(\theta)\|_2^2=1+\sum_{i=2}^N|F_i(\theta)|^2\geq1,
\]
with empty sums when $N=1$.  Nonvanishing is therefore proved before
Lemma~\ref{lem:normalized-curve} is applied to $v=G$.

Proposition~\ref{prop:coordinate-envelope} gives
$|G_i'(x)|\leq D_*$ for all $i$.  Since $D_*\geq0$,
\[
\|G'(x)\|_2^2
=\sum_{i=1}^N|G_i'(x)|^2
\leq\sum_{i=1}^ND_*^2=ND_*^2,
\qquad
\|G'(x)\|_2\leq\sqrt N\,D_*.
\]
Lemma~\ref{lem:normalized-curve} supplies the displayed projector identity,
whose denominator is at least one and whose numerator projector has operator
norm at most one.  Hence
\[
\|\gamma_G'(x)\|_2
\leq\frac{\|G'(x)\|_2}{\|G(x)\|_2}
\leq\sqrt N\,D_*.
\]

Since $F(\theta)=G(x(\theta))$, normalization commutes exactly with
evaluation:
\[
\gamma_F(\theta)
=\frac{G(x(\theta))}{\|G(x(\theta))\|_2}
=\gamma_G(x(\theta)).
\]
The vector chain rule \citep{rudin1976principles} and
$x'(\theta)=h^{-1}$ give
\[
\gamma_F'(\theta)
=\gamma_G'(x(\theta))x'(\theta)
=\frac1h\gamma_G'\!\left(\frac{\theta-c}{h}\right),
\]
including one-sided endpoints by continuity.  Therefore
\[
\|\gamma_F'(\theta)\|_2
\leq\frac{\sqrt N\,D_*}{h},
\]
and taking the defining essential supremum proves the result.

If $N=1$, the anchor forces $F=G=(1)$, so the projector and speed are zero;
also $\Delta=D_*=0$.  If $q=0$, then $M=B_P=0$ and the rate becomes
$\sqrt N\Delta B_Q/h$.  If the projector annihilates $G'$, both normalized
derivatives vanish.  For $G(x)=(1,x/\delta)$ with
$0<\delta\leq1$, one has
\[
B_Q=\frac1\delta,
\qquad
D_*=\frac1\delta,
\qquad
\gamma_G(x)=\frac{(1,x/\delta)}{\sqrt{1+(x/\delta)^2}},
\qquad
\|\gamma_G'(x)\|_2
=\frac{1/\delta}{1+(x/\delta)^2}.
\]
Thus its exact speed supremum is $1/\delta$, while the general certificate is
$\sqrt2/\delta$; both preserve the required scale.
\end{proof}

\subsection{Central incidence geometry and coefficient volume}
\label{app:central-incidence}

\begin{lemma}[Tangential Jacobians of the central incidence]
\label{lem:incidence-jacobians}
Under Assumption~\ref{assump:parameter-regime} and
Proposition~\ref{prop:projective-speed}, let
$J\subseteq\operatorname{int}(\Theta)$ be open, put
$K^\circ=(-R,R)^N$, and write $\gamma=\gamma_F$.  Define
\[
\mathcal S_J
:=\{(\theta,a)\in J\times K^\circ:
       \langle a,\gamma(\theta)\rangle=0\},
\qquad
\pi(\theta,a):=a,
\qquad
\tau(\theta,a):=\theta.
\]
Then $\mathcal S_J$ is an embedded $C^1$, countably $N$-rectifiable
hypersurface.  With
\[
u(\theta,a):=\langle a,\gamma'(\theta)\rangle,
\qquad
n(\theta,a):=\frac{(u(\theta,a),\gamma(\theta))}{\sqrt{1+u(\theta,a)^2}},
\]
$n$ is a unit normal and
\[
J_{\mathcal S_J}\pi
=\frac{|u|}{\sqrt{1+u^2}},
\qquad
J_{\mathcal S_J}\tau
=\frac1{\sqrt{1+u^2}}.
\]
Consequently, including at $u=0$,
\begin{equation}
J_{\mathcal S_J}\pi
=|\langle a,\gamma_F'(\theta)\rangle|
 J_{\mathcal S_J}\tau.
\label{eq:appendix-1}
\end{equation}
\end{lemma}

\begin{proof}
On $U=J\times\mathbb R^N$, define
\[
g(\theta,a):=\langle a,\gamma(\theta)\rangle.
\]
Proposition~\ref{prop:projective-speed} gives $\gamma\in C^1$ and
$\|\gamma\|_2=1$, so
\[
\nabla g(\theta,a)
=(\langle a,\gamma'(\theta)\rangle,\gamma(\theta))=(u,\gamma),
\qquad
\|\nabla g(\theta,a)\|_2^2=1+u^2\geq1.
\]
At each incidence, some coefficient component $\gamma_k(\theta)$ is nonzero,
and $\partial g/\partial a_k=\gamma_k(\theta)$.  The Euclidean $C^1$
implicit-function theorem \citep{rudin1976principles} therefore writes the
zero set locally as a $C^1$ graph over the other $N$ coordinates.  Thus it is
an embedded $C^1$ hypersurface with tangent space $\ker Dg$ and unit normal
$\nabla g/\|\nabla g\|_2$, exactly at the regularity proved above.  No
root-transversality premise is needed: tangent roots $u=0$ remain regular
because the coefficient gradient is the unit vector $\gamma$.

Let $e_0=(1,0,\ldots,0)\in\mathbb R^{N+1}$ be the $\theta$-direction and
$T=T_{(\theta,a)}\mathcal S_J=n^\perp$.  Choose an orthonormal basis
$v_1,\ldots,v_N$ of $T$ and set $s_i=\langle v_i,e_0\rangle$.  Coefficient
projection deletes only the $e_0$ component, so
\[
\langle\pi v_i,\pi v_j\rangle=\delta_{ij}-s_is_j.
\]
The determinant of this rank-one perturbation and the orthonormal basis
$v_1,\ldots,v_N,n$ give
\[
\begin{aligned}
(J_{\mathcal S_J}\pi)^2
&=\det(I_N-ss^{\mathsf T})=1-\|s\|_2^2,\\
1=\|e_0\|_2^2
&=\sum_{i=1}^N\langle e_0,v_i\rangle^2
 +\langle e_0,n\rangle^2
=\|s\|_2^2+\langle e_0,n\rangle^2.
\end{aligned}
\]
Therefore
\[
J_{\mathcal S_J}\pi
=|\langle e_0,n\rangle|
=\frac{|u|}{\sqrt{1+u^2}}.
\]
The tangential gradient of the scalar coordinate $\tau$ is $P_Te_0$, and
hence
\[
J_{\mathcal S_J}\tau
=\|P_Te_0\|_2
=\sqrt{1-\langle e_0,n\rangle^2}
=\frac1{\sqrt{1+u^2}}.
\]
This quantity is strictly positive for every finite $a$.  Multiplication by
$|u|$ proves \eqref{eq:appendix-1} without division by $u$ or selection of a root branch.
\end{proof}

\begin{lemma}[Null boundary and degenerate central classes]
\label{lem:central-null-classes}
Under Assumptions~\ref{assump:parameter-regime} and
\ref{assump:anchored-unit-range}, and
Proposition~\ref{prop:projective-speed}, let $I\subseteq\Theta$ be an
interval and write $\gamma=\gamma_F$.  Then:
\begin{enumerate}
\item For every $a$ and $\theta\in I$,
\begin{equation}
\langle a,F(\theta)\rangle=0
\quad\Longleftrightarrow\quad
\langle a,\gamma(\theta)\rangle=0.
\label{eq:appendix-2}
\end{equation}
\item Coefficients producing a root in
$I\setminus\operatorname{int}(I)$ lie in at most two proper linear
hyperplanes and are Lebesgue-null.
\item If $I\neq\varnothing$, the set
\[
\mathcal Z_I:=\{a:\langle a,F(\theta)\rangle=0
                  \text{ for every }\theta\in I\}
\]
is a proper null linear subspace and has probability zero under every
absolutely continuous law.
\item If $L\subseteq\Theta$ is nondegenerate and $\gamma'=0$ on its
interior, then $\gamma$ is constant on $L$; the root coefficients over $L$
form one proper hyperplane while the sweep integrand vanishes.
\item If $N=1$, then $F=\gamma=(1)$, the root coefficient set is $\{0\}$,
and both sides of Proposition~\ref{prop:central-volume} below are zero for
every interval.
\end{enumerate}
\end{lemma}

\begin{proof}
Proposition~\ref{prop:projective-speed} gives
$F=\|F\|_2\gamma$ with $\|F\|_2>0$, proving \eqref{eq:appendix-2}.

An interval differs from its ordinary interior by at most two endpoints.  At
a fixed endpoint $t$, the root coefficients form
\[
H_t:=\{a:\langle a,\gamma(t)\rangle=0\}=\gamma(t)^\perp.
\]
This hyperplane is proper because $\|\gamma(t)\|_2=1$.  An orthogonal change
of coordinates sends it to $\{a_1=0\}$, which is null by Fubini's theorem
\citep{bogachev2007measure}; orthogonal
maps preserve Lebesgue measure, and a finite union remains null.

If $I\neq\varnothing$, choose $t_0\in I$.  The set $\mathcal Z_I$ is an
intersection of kernels and hence a linear subspace, with
\[
\mathcal Z_I\subseteq\{a:\langle a,F(t_0)\rangle=0\}.
\]
The containing hyperplane is proper because $F_1(t_0)=1$.  Thus
$\mathcal Z_I$ is null, and absolute continuity gives its probability-zero
consequence, including the zero coefficient vector.

If $\gamma'=0$ on the interior of $L$, the scalar mean-value theorem
\citep{rudin1976principles} applied
coordinatewise makes $\gamma$ constant there, and continuity extends that
value to included endpoints.  The root set is then one fixed hyperplane and
$\langle a,\gamma'\rangle=0$ throughout the interval interior.

For $N=1$, the anchor gives $F=\gamma=(1)$.  A root on a nonempty interval
exists only for $a=0$; on an empty interval none exists.  The coefficient
volume is zero.  Also $\gamma'=0$ and
$\gamma(\theta)^\perp\cap[-R,R]=\{0\}$; although
$\mathcal H^0(\{0\})=1$, the section integrand is zero.
\end{proof}

\begin{proposition}[Central root-set coefficient-volume inequality]
\label{prop:central-volume}
Under Assumptions~\ref{assump:parameter-regime} and
\ref{assump:anchored-unit-range}, Proposition~\ref{prop:projective-speed},
and Lemmas~\ref{lem:incidence-jacobians} and
\ref{lem:central-null-classes}, every interval $I\subseteq\Theta$, with any
endpoint convention, satisfies
\begin{equation}
\begin{aligned}
&\operatorname{Leb}^N
\{a\in[-R,R]^N:\exists\theta\in I,
  \ \langle a,F(\theta)\rangle=0\}\\
&\quad\leq
\int_I\int_{\gamma_F(\theta)^\perp\cap[-R,R]^N}
 |\langle\gamma_F'(\theta),a\rangle|
 \,d\mathcal H^{N-1}(a)\,d\theta.
\end{aligned}
\label{eq:appendix-3}
\end{equation}
The bound requires neither simple roots nor transversality.  Tangent
incidences have null coefficient image, distinct roots increase projection
multiplicity, and identically zero combinations may have infinite
multiplicity only on the null subspace in
Lemma~\ref{lem:central-null-classes}.
\end{proposition}

\begin{proof}
Lemma~\ref{lem:central-null-classes} proves the assertion for $N=1$, so
assume $N\geq2$.  Put
\[
K=[-R,R]^N,
\qquad K^\circ=(-R,R)^N,
\qquad J=\operatorname{int}(I),
\qquad \gamma=\gamma_F.
\]
If $J=\varnothing$, then $I$ is empty or a singleton; its root set is empty
or a proper hyperplane, and both the coefficient volume and outer integral
are zero.  Suppose $J\neq\varnothing$.

The root coefficient set is Borel.  Indeed, choose compact intervals
$J_m\subset J$ with $J_m\uparrow J$.  The zero set of the continuous map
$(\theta,a)\mapsto\langle a,\gamma(\theta)\rangle$ in $J_m\times K$ is
compact, so its coefficient projection is compact.  The interior-root set is
the countable union of these projections; adjoining at most two endpoint
hyperplanes preserves Borel measurability.

The cube boundary is Lebesgue-null, and endpoint root classes are null by
Lemma~\ref{lem:central-null-classes}.  The left side of \eqref{eq:appendix-3} therefore has
the same volume as $\pi(\mathcal S_J)$.  For $a\in\mathbb R^N$, define the
extended-valued multiplicity
\[
\mathcal N_J(a)
:=\#\{\theta\in J:(\theta,a)\in\mathcal S_J\}.
\]
The indicator of a nonempty fiber is at most $\mathcal N_J(a)$, even for an
infinite fiber.  The area formula for rectifiable sets
\citep{federer1969gmt} applies because
Lemma~\ref{lem:incidence-jacobians} supplies a countably rectifiable
hypersurface and $\pi$ is $1$-Lipschitz.  Hence, in the extended nonnegative
sense,
\begin{equation}
\begin{aligned}
\operatorname{Leb}^N(\pi(\mathcal S_J))
&=\int_{\mathbb R^N}
  \mathbf1_{\{\mathcal N_J(a)\geq1\}}\,da\\
&\leq\int_{\mathbb R^N}\mathcal N_J(a)\,da\\
&=\int_{\mathcal S_J}
  J_{\mathcal S_J}\pi(\theta,a)\,d\mathcal H^N(\theta,a).
\end{aligned}
\label{eq:appendix-4}
\end{equation}
Using \eqref{eq:appendix-1} and then the coarea formula on the rectifiable hypersurface
\citep{federer1969gmt}, with $\tau(\theta,a)=\theta$, transforms the final
integral into
\begin{equation}
\begin{aligned}
&\int_{\mathcal S_J}
 |\langle a,\gamma'(\theta)\rangle|
 J_{\mathcal S_J}\tau(\theta,a)
 \,d\mathcal H^N(\theta,a)\\
&\quad=
\int_J\int_{\gamma(\theta)^\perp\cap K^\circ}
 |\langle\gamma'(\theta),a\rangle|
 \,d\mathcal H^{N-1}(a)\,d\theta.
\end{aligned}
\label{eq:appendix-5}
\end{equation}
The fiber identification is exact: $a\mapsto(\theta,a)$ is an isometry from
$\gamma(\theta)^\perp\cap K^\circ$ onto the corresponding $\tau$-fiber.
The coarea hypotheses are met because \eqref{eq:appendix-1} gives the measurable
nonnegative integrand and $J_{\mathcal S_J}\tau>0$.

Replacing the open cube section in \eqref{eq:appendix-5} by the closed one changes no
$\mathcal H^{N-1}$ integral.  The cube boundary is the union of $2N$ affine
faces.  A central hyperplane $\gamma(\theta)^\perp$ cannot contain a face:
containment would force its tangential normal components to vanish and its
remaining component times $\pm R$ to be zero, contradicting $R>0$ and
$\|\gamma(\theta)\|_2=1$.  Its intersection with each face has dimension at
most $N-2$ and is therefore $\mathcal H^{N-1}$-null.  Replacing $J$ by $I$
also changes no outer integral because at most two points are added.  Equations
\eqref{eq:appendix-4}--\eqref{eq:appendix-5} prove \eqref{eq:appendix-3}.

To verify the root-type scope, define
\[
\mathcal C_J:=\{(\theta,a)\in\mathcal S_J:
 \langle a,\gamma'(\theta)\rangle=0\}.
\]
On this measurable set, Lemma~\ref{lem:incidence-jacobians} gives
$J_{\mathcal S_J}\pi=0$.  A second area-formula application gives
\[
\int_{\mathbb R^N}
 \#(\mathcal C_J\cap\pi^{-1}(a))\,da
=\int_{\mathcal C_J}J_{\mathcal S_J}\pi\,d\mathcal H^N=0,
\]
so $\operatorname{Leb}^N(\pi(\mathcal C_J))=0$.  At a root,
\[
\frac d{d\theta}\langle a,F(\theta)\rangle
=\|F(\theta)\|_2\langle a,\gamma'(\theta)\rangle,
\]
because the term containing $\langle a,\gamma(\theta)\rangle$ vanishes.
Thus tangent and multiple roots lie in this null image.  Distinct roots add
preimages and only increase $\mathcal N_J$.  An identically zero combination
has infinite multiplicity but lies in the proper null subspace from
Lemma~\ref{lem:central-null-classes}.  Endpoints and stationary pieces were
also handled there.  No finite-root or transversality restriction has been
used.
\end{proof}

\subsection{Cube sections and correlated density conversion}
\label{app:cube-density}

\begin{lemma}[Scaled central cube section]
\label{lem:scaled-cube-section}
Under Assumption~\ref{assump:parameter-regime}, suppose $N\geq2$ and
$v\in\mathbb R^N$ satisfies $\|v\|_2=1$.  Then
\begin{equation}
\mathcal H^{N-1}(v^\perp\cap[-R,R]^N)
\leq\sqrt2(2R)^{N-1}.
\label{eq:appendix-6}
\end{equation}
\end{lemma}

\begin{proof}
Let $Q_N=[-1/2,1/2]^N$ and $H=v^\perp$.  Ball's cube-slicing theorem
\citep{ball1986cube} states, in the Euclidean intrinsic-volume convention,
that every central hyperplane section of this unit-volume cube satisfies
\begin{equation}
\mathcal H^{N-1}(Q_N\cap H)\leq\sqrt2.
\label{eq:appendix-7}
\end{equation}
Its hypotheses hold: $N\geq2$, and the unit vector $v$ makes $H$ a
codimension-one Euclidean linear hyperplane through the origin.  No rotation
of the coefficient law is involved.

Let $D_R(y)=2Ry$.  Since $R>0$,
\[
D_R(Q_N)=[-R,R]^N,
\qquad
D_R(H)=H,
\]
and therefore
\[
D_R(Q_N\cap H)=[-R,R]^N\cap H.
\]
Euclidean Hausdorff measure scales under this similarity by
$(2R)^{N-1}$, directly from its covering definition.  Applying the inverse
dilation supplies the reverse scaling equality, so
\[
\begin{aligned}
\mathcal H^{N-1}(H\cap[-R,R]^N)
&=(2R)^{N-1}\mathcal H^{N-1}(H\cap Q_N)\\
&\leq\sqrt2(2R)^{N-1}.
\end{aligned}
\]
The dilation maps the closed cube and all of its faces exactly.  The result
is orientation-uniform and concerns only central, not translated, sections.
\end{proof}

\begin{proposition}[Central sweep under a capped joint law]
\label{prop:central-sweep}
Under Assumptions~\ref{assump:parameter-regime} and
\ref{assump:cube-density-laws}, Proposition~\ref{prop:central-volume}, and,
when $N\geq2$, Lemma~\ref{lem:scaled-cube-section}, every
$\mu\in\mathcal D_{N,R,\kappa}$ and every interval $I\subseteq\Theta$
satisfy
\begin{equation}
\Pr_{\alpha\sim\mu}
[\exists\theta\in I:\langle\alpha,F(\theta)\rangle=0]
\leq A\sqrt{\frac N2}
\int_I\|\gamma_F'(\theta)\|_2\,d\theta.
\label{eq:appendix-8}
\end{equation}
This includes empty, singleton, open, closed, and half-open intervals.
\end{proposition}

\begin{proof}
Fix the deterministic presentation, then an arbitrary
$\mu\in\mathcal D_{N,R,\kappa}$, and then an arbitrary interval $I$.  Define
the Borel set
\[
\mathcal E_I:=\{a\in[-R,R]^N:\exists\theta\in I,
 \ \langle a,F(\theta)\rangle=0\}.
\]
Its measurability is proved in Proposition~\ref{prop:central-volume}.  If
$f_\mu$ is the full joint density, cube support and
$f_\mu\leq\kappa$ almost everywhere give the only probability conversion:
\begin{equation}
\Pr_{\alpha\sim\mu}
[\exists\theta\in I:\langle\alpha,F(\theta)\rangle=0]
=\int_{\mathcal E_I}f_\mu(a)\,da
\leq\kappa\operatorname{Leb}^N(\mathcal E_I).
\label{eq:appendix-9}
\end{equation}
This uses neither conditional densities nor coordinate independence.  The
nonemptiness of the law class also gives
\[
1=\int_{[-R,R]^N}f_\mu(a)\,da
\leq\kappa(2R)^N=A,
\]
although this fact is not used to relax any later coefficient.

If $N=1$, Proposition~\ref{prop:central-volume} gives
$\operatorname{Leb}^1(\mathcal E_I)=0$: the section is $\{0\}$ and its
integrand vanishes.  Moreover the unit scalar curve has
$\gamma_F'=0$.  Thus \eqref{eq:appendix-8} is exactly $0\leq0$, including all interval
conventions, and Ball's positive-dimensional theorem is not invoked.

Assume $N\geq2$.  Equations \eqref{eq:appendix-3} and \eqref{eq:appendix-9} give
\begin{equation}
\begin{aligned}
&\Pr_{\alpha\sim\mu}
[\exists\theta\in I:\langle\alpha,F(\theta)\rangle=0]\\
&\quad\leq\kappa\int_I
\int_{\gamma_F(\theta)^\perp\cap[-R,R]^N}
|\langle\gamma_F'(\theta),a\rangle|
\,d\mathcal H^{N-1}(a)\,d\theta.
\end{aligned}
\label{eq:appendix-10}
\end{equation}
For $a\in[-R,R]^N$, Cauchy--Schwarz
\citep{rudin1976principles} and cube support give
\begin{equation}
|\langle\gamma_F'(\theta),a\rangle|
\leq\|\gamma_F'(\theta)\|_2\|a\|_2
\leq R\sqrt N\,\|\gamma_F'(\theta)\|_2.
\label{eq:appendix-11}
\end{equation}
Lemma~\ref{lem:scaled-cube-section}, applied to the unit vector
$\gamma_F(\theta)$, therefore yields
\begin{equation}
\begin{aligned}
&\int_{\gamma_F(\theta)^\perp\cap[-R,R]^N}
|\langle\gamma_F'(\theta),a\rangle|
\,d\mathcal H^{N-1}(a)\\
&\quad\leq
R\sqrt N\,\|\gamma_F'(\theta)\|_2
\mathcal H^{N-1}(\gamma_F(\theta)^\perp\cap[-R,R]^N)\\
&\quad\leq
R\sqrt N\,\sqrt2(2R)^{N-1}
\|\gamma_F'(\theta)\|_2.
\end{aligned}
\label{eq:appendix-12}
\end{equation}
Only a central linear section is used.  Substitution into \eqref{eq:appendix-10} gives the
coefficient
\begin{equation}
\begin{aligned}
\kappa R\sqrt N\,\sqrt2(2R)^{N-1}
&=\kappa(2R)^N\frac{\sqrt{2N}}2\\
&=\kappa(2R)^N\sqrt{\frac N2}
=A\sqrt{\frac N2},
\end{aligned}
\label{eq:appendix-13}
\end{equation}
where $R(2R)^{N-1}=(2R)^N/2$.  This is equality, not a term absorption, and
proves \eqref{eq:appendix-8}.  The upper bound need not be at most one and is not clipped.

Endpoint inclusion does not alter the argument.  Proposition
~\ref{prop:central-volume} already accounts for endpoint hyperplanes, tangent
and multiple roots, stationary pieces, infinite fibers, and identically zero
combinations.  A stationary curve makes its incidence integral and event
probability zero.  The dilation in Lemma~\ref{lem:scaled-cube-section}
changes only deterministic section measure; it neither transforms the
coefficient vector nor its density.  Since the law and interval were
arbitrary after the deterministic presentation was fixed, the required
quantifier order follows without a union bound over laws, intervals, roots,
or orientations.
\end{proof}

\subsection{Central rates and ordered suprema}
\label{app:central-rate}

\begin{lemma}[Projective speed controls interval length]
\label{lem:speed-length}
Under Assumption~\ref{assump:parameter-regime} and
Proposition~\ref{prop:projective-speed}, every interval
$I\subseteq\Theta$, with any endpoint convention, satisfies
\begin{equation}
\int_I\|\gamma_F'(\theta)\|_2\,d\theta
\leq\Gamma_{\mathrm{proj}}(F)|I|.
\label{eq:appendix-14}
\end{equation}
\end{lemma}

\begin{proof}
Proposition~\ref{prop:projective-speed} gives
$\gamma_F\in C^1(\Theta;\mathbb R^N)$, so
$\theta\mapsto\|\gamma_F'(\theta)\|_2$ is nonnegative and measurable.  By
the definition of the essential supremum, outside a null subset of $\Theta$,
\[
\|\gamma_F'(\theta)\|_2\leq\Gamma_{\mathrm{proj}}(F).
\]
Integration over $I$ gives
\[
\int_I\|\gamma_F'(\theta)\|_2\,d\theta
\leq\int_I\Gamma_{\mathrm{proj}}(F)\,d\theta
=\Gamma_{\mathrm{proj}}(F)|I|.
\]
Changing endpoint inclusion affects neither the integral nor the length, so
the assertion also includes empty and singleton intervals.
\end{proof}

\begin{proposition}[Central rate bridge and supremum closure]
\label{prop:central-rate}
Under Assumptions~\ref{assump:parameter-regime},
\ref{assump:balcan-common-chain},
\ref{assump:anchored-unit-range}, and
\ref{assump:cube-density-laws}, Propositions~\ref{prop:projective-speed} and
\ref{prop:central-sweep}, and Lemma~\ref{lem:speed-length}, every
$\mu\in\mathcal D_{N,R,\kappa}$ and every positive-length interval
$I\subseteq\Theta$ satisfy
\begin{equation}
\begin{aligned}
&\Pr_{\alpha\sim\mu}
[\exists\theta\in I:\langle\alpha,F(\theta)\rangle=0]\\
&\quad\leq A\sqrt{\frac N2}\,\Gamma_{\mathrm{proj}}(F)|I|\\
&\quad\leq
\frac{AN\Delta B_Q(1+qB_P)}{\sqrt2\,h}|I|.
\end{aligned}
\label{eq:appendix-15}
\end{equation}
Moreover,
\begin{equation}
C^{\mathrm{Pf}}_{\mathcal D}(F;\Theta)
\leq A\sqrt{\frac N2}\,\Gamma_{\mathrm{proj}}(F)
\leq\frac{AN\Delta B_Q(1+qB_P)}{\sqrt2\,h}.
\label{eq:appendix-16}
\end{equation}
The last coefficient is literally
\begin{equation}
\frac{(2R)^N\kappa N\Delta B_Q(1+qB_P)}{\sqrt2\,h},
\label{eq:appendix-17}
\end{equation}
has degree-zero dependence on $M$, reduces to
$AN\Delta B_Q/(\sqrt2 h)$ when $q=0$, and becomes
$\sqrt2/\delta$ when
$N=2$, $A=h=\Delta=1$, $q=B_P=0$, and $B_Q=1/\delta$.
\end{proposition}

\begin{proof}
Fix the deterministic presentation and full parameter tuple first.  Then fix
an arbitrary $\mu\in\mathcal D_{N,R,\kappa}$ and an arbitrary interval $I$,
initially allowing zero length.  Define
\[
\mathsf Z_I:=\{\alpha\in\mathbb R^N:\exists\theta\in I,
 \ \langle\alpha,F(\theta)\rangle=0\}.
\]
Proposition~\ref{prop:central-sweep} and
Lemma~\ref{lem:speed-length} give
\begin{equation}
\begin{aligned}
\Pr_{\alpha\sim\mu}(\mathsf Z_I)
&\leq A\sqrt{\frac N2}
 \int_I\|\gamma_F'(\theta)\|_2\,d\theta\\
&\leq A\sqrt{\frac N2}\,
 \Gamma_{\mathrm{proj}}(F)|I|.
\end{aligned}
\label{eq:appendix-18}
\end{equation}
Proposition~\ref{prop:projective-speed} applies to the same deterministic
feature and Euclidean norm.  Multiplying its bound by the nonnegative factor
$A\sqrt{N/2}|I|$ gives
\begin{equation}
\begin{aligned}
A\sqrt{\frac N2}\,\Gamma_{\mathrm{proj}}(F)|I|
&\leq A\sqrt{\frac N2}
 \frac{\sqrt N\,\Delta B_Q(1+qB_P)}h|I|\\
&=\frac{AN\Delta B_Q(1+qB_P)}{\sqrt2\,h}|I|,
\end{aligned}
\label{eq:appendix-19}
\end{equation}
where the complete simplification is
\begin{equation}
\sqrt{\frac N2}\sqrt N
=\sqrt{\frac{N^2}{2}}
=\frac N{\sqrt2}
\qquad(N\geq1).
\label{eq:appendix-20}
\end{equation}
No term is dropped.  Equations \eqref{eq:appendix-18}--\eqref{eq:appendix-20} prove \eqref{eq:appendix-15}.  If $I$ is
empty or a singleton, the speed integral and $|I|$ are zero, so the event has
probability zero before any quotient.  Endpoint variants have the same
length and integral, while Proposition~\ref{prop:central-sweep} retains the
root-event boundary handling.

Now restrict to $|I|>0$ and divide \eqref{eq:appendix-15} by this positive length:
\begin{equation}
\frac{\Pr_{\alpha\sim\mu}(\mathsf Z_I)}{|I|}
\leq A\sqrt{\frac N2}\,\Gamma_{\mathrm{proj}}(F)
\leq\frac{AN\Delta B_Q(1+qB_P)}{\sqrt2\,h}.
\label{eq:appendix-21}
\end{equation}
For fixed $\mu$, first take the interval supremum.  The two right sides are
independent of $I$, so
\begin{equation}
\sup_{\substack{I\subseteq\Theta\ \mathrm{interval}\\|I|>0}}
\frac{\Pr_{\alpha\sim\mu}(\mathsf Z_I)}{|I|}
\leq A\sqrt{\frac N2}\,\Gamma_{\mathrm{proj}}(F)
\leq\frac{AN\Delta B_Q(1+qB_P)}{\sqrt2\,h}.
\label{eq:appendix-22}
\end{equation}
Only then take the law supremum:
\begin{equation}
\begin{aligned}
&\sup_{\mu\in\mathcal D_{N,R,\kappa}}
 \sup_{\substack{I\subseteq\Theta\ \mathrm{interval}\\|I|>0}}
 \frac{\Pr_{\alpha\sim\mu}(\mathsf Z_I)}{|I|}\\
&\quad\leq A\sqrt{\frac N2}\,\Gamma_{\mathrm{proj}}(F)
\leq\frac{AN\Delta B_Q(1+qB_P)}{\sqrt2\,h}.
\end{aligned}
\label{eq:appendix-23}
\end{equation}
The left side is the definition of
$C^{\mathrm{Pf}}_{\mathcal D}(F;\Theta)$, proving \eqref{eq:appendix-16}.

Substituting $A=(2R)^N\kappa$ gives \eqref{eq:appendix-17} by equality.  The parameter $M$
does not occur, so its dependence is degree zero for fixed $B_P$.  At $q=0$,
$M=B_P=0$ and $1+qB_P=1$.  For
$N=2$, $A=h=\Delta=1$, $q=B_P=0$, and $B_Q=1/\delta$, the coefficient is
\begin{equation}
\frac{1\cdot2\cdot1\cdot(1/\delta)\cdot1}{\sqrt2\cdot1}
=\frac{\sqrt2}{\delta}.
\label{eq:appendix-24}
\end{equation}
If $\Gamma_{\mathrm{proj}}(F)=0$, \eqref{eq:appendix-18} gives zero probability for every
law and interval, hence zero capacity; no division by speed occurs.  When
$N=1$, the anchor gives $F=(1)$, $\Delta=0$, and all displayed chains reduce
to $0\leq0$.  The only division in the closure is by the explicitly positive
interval length.
\end{proof}

\subsection{Finite affine charts and measurable exhaustion}
\label{app:affine-charts}

\begin{lemma}[Finite exhausted-chart legality]
\label{lem:finite-chart-legality}
Under Assumptions~\ref{assump:parameter-regime},
\ref{assump:balcan-common-chain},
\ref{assump:anchored-unit-range}, and
\ref{assump:affine-chart-data}, and
Proposition~\ref{prop:coordinate-envelope}, fix an interval
$I\subseteq\Theta$, an index $1\leq j\leq N$, and $m\geq1$.  Put
\[
E_{j,m}:=\{\theta\in E_j:|F_j(\theta)|\geq1/m\}.
\]
Index $\beta\in[-R,R]^{N-1}$ by $i\neq j$, using the zero-dimensional
singleton convention at $N=1$, and define
\[
\Psi_j(\theta,\beta)_i=
\begin{cases}
\beta_i,&i\neq j,\\
T_j(\theta,\beta),&i=j.
\end{cases}
\]
Then $T_j$ and $\Psi_j$ are measurable and Lipschitz on
$E_{j,m}\times[-R,R]^{N-1}$.  Their restrictions to measurable subsets
satisfy the regularity hypotheses of the equal-dimensional area formula.  No
Lipschitz constant uniform in $m$ is asserted.
\end{lemma}

\begin{proof}
For every $\theta\in\Theta$, the unit-box coefficient estimate gives
\begin{equation}
|F_i(\theta)|
=|Q_i(x(\theta),\eta(x(\theta)))|
\leq\|\operatorname{coeff}(Q_i)\|_1
\leq B_Q.
\label{eq:appendix-25}
\end{equation}
Proposition~\ref{prop:coordinate-envelope} and the affine coordinate change
give
\begin{equation}
|F_i'(\theta)|\leq L_F:=D_*/h.
\label{eq:appendix-26}
\end{equation}
The proof-local quantities
\begin{equation}
B_H:=\|F_0\|_\infty+(N-1)RB_Q,
\qquad
L_H:=\|F_0'\|_\infty+(N-1)RL_F
\label{eq:appendix-27}
\end{equation}
are finite because $F_0\in C^1(\Theta)$, $\Theta$ is compact, and the
remaining terms are setting quantities or the bound \eqref{eq:appendix-26}.  Define
\[
H_j(\theta,\beta)
:=F_0(\theta)+\sum_{i\neq j}\beta_iF_i(\theta),
\qquad
T_j=-H_j/F_j.
\]
Then
\[
|H_j(\theta,\beta)|\leq B_H,
\]
and, for $\theta,s\in\Theta$ and
$\beta,\zeta\in[-R,R]^{N-1}$,
\begin{equation}
|H_j(\theta,\beta)-H_j(s,\zeta)|
\leq L_H|\theta-s|+B_Q\|\beta-\zeta\|_1.
\label{eq:appendix-28}
\end{equation}
If $\theta,s\in E_{j,m}$, both denominators have magnitude at least $1/m$.
The quotient identity, \eqref{eq:appendix-26}, and \eqref{eq:appendix-28} yield
\begin{equation}
\begin{aligned}
|T_j(\theta,\beta)-T_j(s,\zeta)|
&\leq m|H_j(\theta,\beta)-H_j(s,\zeta)|\\
&\quad+m^2|H_j(s,\zeta)|\,|F_j(\theta)-F_j(s)|\\
&\leq(mL_H+m^2B_HL_F)|\theta-s|
 +mB_Q\|\beta-\zeta\|_1.
\end{aligned}
\label{eq:appendix-29}
\end{equation}
Thus $T_j$ is Lipschitz even if the measurable finite-level set is
disconnected; inserting the unchanged beta coordinates makes $\Psi_j$
Lipschitz.  At interior points,
\begin{equation}
|\partial_\theta T_j(\theta,\beta)|
\leq mL_H+m^2B_HL_F,
\qquad
|\partial_{\beta_i}T_j(\theta,\beta)|
\leq mB_Q.
\label{eq:appendix-30}
\end{equation}
The same bounds hold one-sided at ambient endpoints, and the finite-level
Jacobian is bounded on its finite product domain.

On the open nonzero-pivot locus in $\operatorname{int}\Theta$, the quotient
formulas are $C^1$.  Hence at almost every density point of a measurable
restricted domain, the approximate derivative of the Lipschitz restriction
agrees with this $C^1$ derivative.  The ambient endpoints are null affine
slices, as established in Lemma~\ref{lem:affine-null-classes} below.
Measurability follows from measurability of $E_j$ and continuity of the
quotient on the nonzero-pivot locus.  When $N=1$, all beta sums and norms are
empty and the argument is the same scalar quotient estimate.
\end{proof}

\begin{lemma}[Exact affine coefficient chart and Jacobian]
\label{lem:affine-chart-jacobian}
Under the assumptions and fixed objects of
Lemma~\ref{lem:finite-chart-legality}, define
\[
D_{j,m}:=\{(\theta,\beta):\theta\in E_{j,m},
 \ \beta\in[-R,R]^{N-1},\ |T_j(\theta,\beta)|\leq R\}.
\]
Then $D_{j,m}$ is measurable,
$\Psi_j(D_{j,m})\subseteq[-R,R]^N$, and every point in $D_{j,m}$ satisfies
\begin{equation}
F_0(\theta)+\langle\Psi_j(\theta,\beta),F(\theta)\rangle=0.
\label{eq:appendix-31}
\end{equation}
Conversely, if $\alpha\in[-R,R]^N$ has a root at
$\theta\in E_{j,m}$, then, with $\beta=\alpha_{-j}$,
$(\theta,\beta)\in D_{j,m}$ and $\Psi_j(\theta,\beta)=\alpha$.  At every
interior chart point,
\begin{equation}
|\det D\Psi_j(\theta,\beta)|
=|\partial_\theta T_j(\theta,\beta)|.
\label{eq:appendix-32}
\end{equation}
At a represented root,
\begin{equation}
\frac d{d\theta}
\bigl(F_0(\theta)+\langle\alpha,F(\theta)\rangle\bigr)
=-F_j(\theta)\partial_\theta T_j(\theta,\beta),
\label{eq:appendix-33}
\end{equation}
so tangent roots have zero chart Jacobian.
\end{lemma}

\begin{proof}
Measurability follows from Lemma~\ref{lem:finite-chart-legality}.  The domain
definition places every $\beta_i$ and the inserted coordinate $T_j$ in
$[-R,R]$, proving cube membership.  By the chart definition,
\[
\begin{aligned}
F_0(\theta)+\langle\Psi_j(\theta,\beta),F(\theta)\rangle
&=F_0(\theta)+\sum_{i\neq j}\beta_iF_i(\theta)
 +T_j(\theta,\beta)F_j(\theta)\\
&=H_j(\theta,\beta)-H_j(\theta,\beta)=0.
\end{aligned}
\]
Conversely, if the affine expression vanishes and $F_j(\theta)\neq0$, solving
for $\alpha_j$ gives $\alpha_j=T_j(\theta,\alpha_{-j})$.  Cube membership of
$\alpha$ supplies $\alpha_{-j}\in[-R,R]^{N-1}$ and $|T_j|\leq R$, proving
the converse.

Permute output coordinates so coordinate $j$ is first, and retain domain
coordinates $(\theta,(\beta_i)_{i\neq j})$.  The derivative matrix is
\[
\begin{pmatrix}
\partial_\theta T_j&(\partial_{\beta_i}T_j)_{i\neq j}\\
0&I_{N-1}
\end{pmatrix}.
\]
The coordinate permutation has determinant of magnitude one, proving \eqref{eq:appendix-32}.
For $N=1$, the lower-right block is the $0\times0$ identity with determinant
one, and the formula becomes $|\Psi_1'|=|T_1'|$.

Differentiate
$H_j(\theta,\beta)+T_j(\theta,\beta)F_j(\theta)=0$ with beta fixed.  At
$\alpha=\Psi_j(\theta,\beta)$ this gives
\[
F_0'(\theta)+\sum_{i\neq j}\alpha_iF_i'(\theta)
 +\alpha_jF_j'(\theta)
=-F_j(\theta)\partial_\theta T_j(\theta,\beta),
\]
which is \eqref{eq:appendix-33}.  Since the pivot is nonzero, tangency is equivalent to a
zero chart derivative.
\end{proof}

\begin{lemma}[Null affine endpoint and degeneracy classes]
\label{lem:affine-null-classes}
Under Assumptions~\ref{assump:anchored-unit-range},
\ref{assump:cube-density-laws}, and
\ref{assump:affine-chart-data}, let $I\subseteq\Theta$ be nonempty.  For
every fixed $\theta_0\in I$,
\[
\mathcal H_{\theta_0}
:=\{\alpha\in\mathbb R^N:
F_0(\theta_0)+\langle\alpha,F(\theta_0)\rangle=0\}
\]
is a proper affine hyperplane and has probability zero under every
$\mu\in\mathcal D_{N,R,\kappa}$.  Hence coefficients having only an included
endpoint root are law-null.  The set
\[
\mathcal Z_I:=\{\alpha:F_0(\theta)+\langle\alpha,F(\theta)\rangle=0
                 \text{ for every }\theta\in I\}
\]
is empty or a proper affine subspace and is also law-null.
\end{lemma}

\begin{proof}
The anchor gives $F_1(\theta_0)=1$, so
\[
\alpha_1=-F_0(\theta_0)-\sum_{i=2}^N\alpha_iF_i(\theta_0)
\]
is a proper affine graph.  For each fixed
$(\alpha_2,\ldots,\alpha_N)$, its fiber in the first coordinate is a single
point and is one-dimensional Lebesgue-null.  Tonelli's theorem
\citep{bogachev2007measure} gives
$\operatorname{Leb}^N(\mathcal H_{\theta_0})=0$; at $N=1$, the empty-sum
graph is itself one point.  Absolute continuity gives law nullity.

There are at most two included endpoints, so their root classes lie in a
finite union of such hyperplanes.  Moreover
$\mathcal Z_I=\bigcap_{\theta\in I}\mathcal H_\theta$ is empty or affine,
and for any $\theta_0\in I$ it lies in the proper null hyperplane
$\mathcal H_{\theta_0}$.  If $I$ is empty, the root event is empty and no
degenerate exception is needed.  Only the full joint absolute continuity has
been used.
\end{proof}

\begin{proposition}[Finite-level chart area and density bound]
\label{prop:finite-chart-area}
Under Assumptions~\ref{assump:parameter-regime},
\ref{assump:balcan-common-chain},
\ref{assump:anchored-unit-range},
\ref{assump:cube-density-laws}, and
\ref{assump:affine-chart-data}, Proposition~\ref{prop:coordinate-envelope},
and Lemmas~\ref{lem:finite-chart-legality},
\ref{lem:affine-chart-jacobian}, and
\ref{lem:affine-null-classes}, fix an interval $I$, $m\geq1$, and
$\mu\in\mathcal D_{N,R,\kappa}$.  Define
\[
\mathcal B_{j,m}:=\Psi_j(D_{j,m}),
\qquad
\mathcal R_m:=\bigcup_{j=1}^N\mathcal B_{j,m}.
\]
These coefficient sets are measurable in their Lebesgue and $\mu$
completions, and
\begin{equation}
\begin{aligned}
\Pr_{\alpha\sim\mu}[\alpha\in\mathcal R_m]
&\leq\kappa\sum_{j=1}^N
\int_{E_{j,m}}\int_{[-R,R]^{N-1}}
\mathbf1_{\{|T_j(\theta,\beta)|\leq R\}}
|\partial_\theta T_j(\theta,\beta)|
\,d\beta\,d\theta\\
&\leq\kappa\sum_{j=1}^N
\int_{E_{j,m}}\int_{[-R,R]^{N-1}}
|\partial_\theta T_j(\theta,\beta)|
\,d\beta\,d\theta.
\end{aligned}
\label{eq:appendix-34}
\end{equation}
Both finite-level right sides are finite.  Injectivity and a bound on root
multiplicity are unnecessary.
\end{proposition}

\begin{proof}
If $I$ is empty, every set and integral in \eqref{eq:appendix-34} is empty, so assume $I$ is
nonempty.  Fix $j$.  Lemma~\ref{lem:finite-chart-legality} makes $\Psi_j$
Lipschitz on the finite-level domain containing $D_{j,m}$.  Away from the at
most two endpoints of $\Theta$, its approximate derivative agrees almost
everywhere with the $C^1$ chart derivative.  Endpoint slices map into the
null hyperplanes in Lemma~\ref{lem:affine-null-classes}.  The
equal-dimensional area formula \citep{federer1969gmt} and
Lemma~\ref{lem:affine-chart-jacobian} therefore give
\begin{equation}
\begin{aligned}
\operatorname{Leb}^N(\mathcal B_{j,m})
&\leq\int_{D_{j,m}}
 |\det D\Psi_j(\theta,\beta)|\,d\beta\,d\theta\\
&=\int_{E_{j,m}}\int_{[-R,R]^{N-1}}
 \mathbf1_{\{|T_j(\theta,\beta)|\leq R\}}
 |\partial_\theta T_j(\theta,\beta)|
 \,d\beta\,d\theta.
\end{aligned}
\label{eq:appendix-35}
\end{equation}
The formula applies to a Lipschitz extension of the map from its measurable
domain.  If that domain is not Borel, write it as a Borel set plus a null set.
The equal-dimensional Lipschitz image of the null part is null, while the
Borel image is analytic and hence universally measurable
\citep{bogachev2007measure}.  This proves the completed-measure assertion
without strengthening the assumed measurability of $E_j$.

The area formula integrates the full preimage-counting function.  Multiple
roots can only increase the right side of \eqref{eq:appendix-35}.  By
Lemma~\ref{lem:affine-chart-jacobian}, tangent roots lie in the zero-Jacobian
part, whose image is null under the same formula.  Affine-identically-zero
coefficients are independently law-null by
Lemma~\ref{lem:affine-null-classes}; neither root simplicity nor
transversality is used.

Every $\mathcal B_{j,m}$ lies in $[-R,R]^N$.  Finite subadditivity and the
single full joint-density cap give
\begin{equation}
\mu(\mathcal R_m)
\leq\kappa\operatorname{Leb}^N(\mathcal R_m)
\leq\kappa\sum_{j=1}^N\operatorname{Leb}^N(\mathcal B_{j,m}).
\label{eq:appendix-36}
\end{equation}
Combining \eqref{eq:appendix-35}--\eqref{eq:appendix-36} proves the first line of \eqref{eq:appendix-34}; removing the
indicator proves the second by nonnegativity.  The finite-level derivative
bounds \eqref{eq:appendix-30} and the finite measure of
$I\times[-R,R]^{N-1}$ make both right sides finite.  At $N=1$, integration
over $[-R,R]^0$ has mass one and the same calculation is one-dimensional.
\end{proof}

\begin{lemma}[Least-level pivot exhaustion]
\label{lem:affine-exhaustion}
Under Assumptions~\ref{assump:anchored-unit-range} and
\ref{assump:affine-chart-data}, and
Lemma~\ref{lem:affine-chart-jacobian}, define
\[
\mathcal R_I:=\{\alpha\in[-R,R]^N:\exists\theta\in I,
 \ F_0(\theta)+\langle\alpha,F(\theta)\rangle=0\}.
\]
With $\mathcal R_m$ as in Proposition~\ref{prop:finite-chart-area}, let
\[
C_{j,m}:=E_{j,m}\setminus E_{j,m-1},
\qquad E_{j,0}:=\varnothing.
\]
Then
\[
E_{j,m}=\bigsqcup_{\ell=1}^m C_{j,\ell},
\]
the cells are measurable, and every $\theta\in E_j$ has a unique least
finite activation level.  If
\[
\widehat D_{j,m}:=D_{j,m}\cap
 (C_{j,m}\times[-R,R]^{N-1}),
\]
then
\begin{equation}
D_{j,m}=\bigsqcup_{\ell=1}^m\widehat D_{j,\ell},
\qquad
\mathcal R_m\subseteq\mathcal R_{m+1},
\qquad
\mathcal R_I=\bigcup_{m=1}^\infty\mathcal R_m.
\label{eq:appendix-37}
\end{equation}
The equality includes endpoint, tangent, multiple, and identically-zero
roots.
\end{lemma}

\begin{proof}
Since $1/(m+1)\leq1/m$, the measurable sets satisfy
$E_{j,m}\subseteq E_{j,m+1}$.  The pairwise disjoint differences telescope
to $E_{j,m}=\bigsqcup_{\ell=1}^m C_{j,\ell}$.  Intersecting with the
measurable restriction $|T_j|\leq R$ gives the stated disjoint identity for
$D_{j,m}$.

For $\theta\in E_j$, the selected pivot is nonzero.  Therefore
\[
\{m\geq1:|F_j(\theta)|\geq1/m\}
\]
is a nonempty terminal subset of the positive integers and has the unique
least element
\begin{equation}
m_*(\theta):=\min\{m\geq1:|F_j(\theta)|\geq1/m\}.
\label{eq:appendix-38}
\end{equation}
Equivalently, $\theta\in C_{j,m_*(\theta)}$.  The partition
$(E_j)_{j=1}^N$ also assigns a unique chart owner, so $(j,m_*(\theta))$ is
the unique least finite chart label.  No lower bound uniform in $\theta$ or
$m$ has been introduced.

The set inclusions imply $D_{j,m}\subseteq D_{j,m+1}$, and the same
$\Psi_j$ is used at every level, hence
$\mathcal R_m\subseteq\mathcal R_{m+1}$.  Each chart image consists of root
coefficients by Lemma~\ref{lem:affine-chart-jacobian}, proving
$\bigcup_m\mathcal R_m\subseteq\mathcal R_I$.

Conversely, let $\alpha\in\mathcal R_I$ have a witnessing root $\theta$.
The partition assigns $\theta$ to a unique $E_j$, and its pivot is nonzero.
At the least level \eqref{eq:appendix-38}, put $\beta=\alpha_{-j}$.  The converse part of
Lemma~\ref{lem:affine-chart-jacobian} gives
$(\theta,\beta)\in D_{j,m_*(\theta)}$ and
$\Psi_j(\theta,\beta)=\alpha$.  Thus
$\alpha\in\mathcal R_{m_*(\theta)}$, proving \eqref{eq:appendix-37}.

For measurability, $\mathcal R_I$ is the coefficient projection of the zero
set of a continuous function on the Borel set
$I\times[-R,R]^N$; it is analytic and universally measurable
\citep{bogachev2007measure}.  The finite chart events have the completed
measurability from Proposition~\ref{prop:finite-chart-area}.  Endpoint and
degenerate classes are not removed from \eqref{eq:appendix-37}; their nullity is a separate
measure statement.  If $I$ is empty, both sides are empty.
\end{proof}

\begin{proposition}[General affine chart inequality]
\label{prop:affine-bound}
Under Assumptions~\ref{assump:parameter-regime},
\ref{assump:balcan-common-chain},
\ref{assump:anchored-unit-range},
\ref{assump:cube-density-laws}, and
\ref{assump:affine-chart-data}, Proposition~\ref{prop:coordinate-envelope},
Proposition~\ref{prop:finite-chart-area}, and
Lemma~\ref{lem:affine-exhaustion}, every
$\mu\in\mathcal D_{N,R,\kappa}$ and every interval $I\subseteq\Theta$
satisfy
\begin{equation}
\begin{aligned}
&\Pr_{\alpha\sim\mu}
[\exists\theta\in I:F_0(\theta)+\langle\alpha,F(\theta)\rangle=0]\\
&\quad\leq\kappa\sum_{j=1}^N
\int_{E_j}\int_{[-R,R]^{N-1}}
\mathbf1_{\{|T_j(\theta,\beta)|\leq R\}}
|\partial_\theta T_j(\theta,\beta)|\,d\beta\,d\theta\\
&\quad\leq\kappa\sum_{j=1}^N
\int_{E_j}\int_{[-R,R]^{N-1}}
|\partial_\theta T_j(\theta,\beta)|\,d\beta\,d\theta.
\end{aligned}
\label{eq:appendix-39}
\end{equation}
Both inequalities hold in $[0,+\infty]$.
\end{proposition}

\begin{proof}
The set equality in Lemma~\ref{lem:affine-exhaustion} and continuity from
below \citep{bogachev2007measure} give
\begin{equation}
\mu(\mathcal R_I)=\lim_{m\to\infty}\mu(\mathcal R_m).
\label{eq:appendix-40}
\end{equation}
Because every admissible law is supported on the coefficient cube, the left
side is precisely the affine root-event probability.

For each $j$, the nonnegative measurable functions
\[
\mathbf1_{E_{j,m}}(\theta)
\mathbf1_{\{|T_j(\theta,\beta)|\leq R\}}
|\partial_\theta T_j(\theta,\beta)|
\]
increase pointwise on $E_j\times[-R,R]^{N-1}$ to the same expression with
$E_j$ in place of $E_{j,m}$.  Apply Proposition
~\ref{prop:finite-chart-area}, then monotone convergence
\citep{bogachev2007measure} and finiteness of the chart index set:
\begin{equation}
\begin{aligned}
\mu(\mathcal R_I)
&\leq\kappa\lim_{m\to\infty}\sum_{j=1}^N
\int_{E_{j,m}}\int_{[-R,R]^{N-1}}
\mathbf1_{\{|T_j|\leq R\}}|\partial_\theta T_j|
\,d\beta\,d\theta\\
&=\kappa\sum_{j=1}^N
\int_{E_j}\int_{[-R,R]^{N-1}}
\mathbf1_{\{|T_j|\leq R\}}|\partial_\theta T_j|
\,d\beta\,d\theta.
\end{aligned}
\label{eq:appendix-41}
\end{equation}
All limits and integrals are in $[0,+\infty]$, so no finiteness of the
limiting chart integral is required.  If it diverges, the right side is
$+\infty$.  Dropping the indicator by nonnegativity proves the second
inequality in \eqref{eq:appendix-39}.  The density cap was used only on the full coefficient
set in Proposition~\ref{prop:finite-chart-area}; no coordinate conditioning
was introduced.  The preceding results include empty and singleton
intervals, endpoints, $N=1$, near-zero pivots, tangent and multiple roots,
infinite fibers, and identically-zero affine combinations.
\end{proof}

\subsection{Exact monic presentation and chart velocities}
\label{app:monic-charts}

\begin{proposition}[Exact normalized monic presentation]
\label{prop:monic-presentation}
Under Assumptions~\ref{assump:parameter-regime},
\ref{assump:balcan-common-chain}, and
\ref{assump:anchored-unit-range}, let $d\geq1$ and let
$J\subset\mathbb R$ be a bounded interval.  For the specialization with
$N=d$, there is a nondegenerate interval
$\Theta=[c-h,c+h]\supseteq J$ for which
\begin{equation}
F_0(\theta)=\theta^d,
\qquad
F_{k+1}(\theta)=\theta^k\quad(0\leq k\leq d-1),
\label{eq:appendix-42}
\end{equation}
has normalized-coordinate representatives
\begin{equation}
Q_0(x)=(c+hx)^d,
\qquad
Q_{k+1}(x)=(c+hx)^k\quad(0\leq k\leq d-1).
\label{eq:appendix-43}
\end{equation}
For every $\theta\in\Theta$ and
$\alpha=(\alpha_0,\ldots,\alpha_{d-1})\in\mathbb R^d$,
\begin{equation}
F_0(\theta)+\langle\alpha,F(\theta)\rangle
=\theta^d+\sum_{k=0}^{d-1}\alpha_k\theta^k
=p_\alpha(\theta).
\label{eq:appendix-44}
\end{equation}
The specialization tuple is
\begin{equation}
q=0,
\qquad M=0,
\qquad B_P=0,
\qquad N=d,
\qquad A=(2R)^d\kappa,
\qquad \Delta_{\mathrm{aug}}=d.
\label{eq:appendix-45}
\end{equation}
The leading coefficient $1$ is deterministic and is not a coordinate of the
random lower-coefficient vector.
\end{proposition}

\begin{proof}
Since $J$ is bounded, choose finite $L<U$ with $J\subseteq[L,U]$ and put
\[
c:=\frac{L+U}{2},
\qquad
h:=\frac{U-L}{2}>0.
\]
Then $\Theta=[L,U]$ contains $J$.  Empty and singleton intervals are included
by choosing a strict margin around them.

Set $q=0$.  The chain is empty, so the conventions give
$M=B_P=0$ and $G_i=Q_i$.  For
$x(\theta)=(\theta-c)/h$,
\[
c+hx(\theta)=c+h\frac{\theta-c}{h}=\theta.
\]
This proves \eqref{eq:appendix-42}--\eqref{eq:appendix-43}.  In particular
$Q_1(x)=(c+hx)^0=1$, so the anchor holds and the chain-range clause is
vacuous.  All outputs are polynomials, hence $F_0\in C^1(\Theta)$, and their
monomial coefficient budgets are finite static presentation data.

There are exactly $d$ random-feature outputs, indexed by
$k=0,\ldots,d-1$, so $N=d$ and
\[
A=(2R)^N\kappa=(2R)^d\kappa.
\]
Because $h>0$, the coefficient of $x^d$ in $(c+hx)^d$ is
$h^d\neq0$, so $\deg Q_0=d$, while
$\deg Q_{k+1}=k\leq d-1$.  The same affine coordinate change has nonzero
slope and preserves these degrees.  Hence
\[
\Delta_{\mathrm{aug}}
:=\max\{\deg F_0,\deg F_1,\ldots,\deg F_d\}=d.
\]
Substitution gives \eqref{eq:appendix-44}.  The leading coefficient occurs only in the
deterministic offset $F_0$, leaving exactly the $d$ lower coefficients in
$\alpha$.
\end{proof}

\begin{lemma}[Prescribed monic pivot partition]
\label{lem:monic-pivot-partition}
Under Assumption~\ref{assump:parameter-regime} and
Proposition~\ref{prop:monic-presentation}, if $d\geq2$, then
\begin{equation}
E_1=J\cap\{|\theta|\leq1\},
\qquad
E_d=J\cap\{|\theta|>1\},
\qquad
E_j=\varnothing\quad(j\notin\{1,d\})
\label{eq:appendix-46}
\end{equation}
is a measurable partition of $J$, with $F_1\neq0$ on $E_1$ and
$F_d\neq0$ on $E_d$.  If $d=1$, the choice $E_1=J$ is measurable and has
$F_1\neq0$.  The conclusions include $\theta=0$, $|\theta|=1$, and empty
cells.
\end{lemma}

\begin{proof}
Intervals and the two sets $\{|\theta|\leq1\}$ and
$\{|\theta|>1\}$ are Borel.  The latter sets are disjoint and cover
$\mathbb R$, proving the partition and its measurability.  Empty-cell pivot
conditions are vacuous.

Proposition~\ref{prop:monic-presentation} gives
$F_1(\theta)=\theta^0=1$.  Thus $0$ and the transition points $\pm1$, when
present, lie in $E_1$ and have pivot one.  For $d\geq2$,
\[
F_d(\theta)=\theta^{d-1},
\]
which is nonzero, indeed has magnitude greater than one, on $E_d$.  At
$d=1$, the sole cell $E_1=J$ has the same constant pivot.  The construction
does not depend on the location of $J$; an unused regime simply gives an
empty cell.
\end{proof}

\begin{lemma}[Constant-pivot chart]
\label{lem:monic-low-chart}
Under Assumption~\ref{assump:parameter-regime},
Proposition~\ref{prop:monic-presentation}, and
Lemma~\ref{lem:monic-pivot-partition}, if $d\geq2$,
$\theta\in E_1$, and
\[
\beta=(\beta_1,\ldots,\beta_{d-1})\in[-R,R]^{d-1},
\]
then
\begin{equation}
T_1(\theta,\beta)
=-\theta^d-\sum_{k=1}^{d-1}\beta_k\theta^k,
\label{eq:appendix-47}
\end{equation}
and
\begin{equation}
|\partial_\theta T_1(\theta,\beta)|
\leq d+R\sum_{k=1}^{d-1}k
=d+\frac{Rd(d-1)}2.
\label{eq:appendix-48}
\end{equation}
If $d=1$, then $T_1(\theta)=-\theta$ and
$|T_1'(\theta)|=1$.
\end{lemma}

\begin{proof}
For $d\geq2$, index nonpivot coordinates by polynomial exponent.  On $E_1$,
the root equation with $\alpha_0$ as pivot is
\[
0=\theta^d+\alpha_0+\sum_{k=1}^{d-1}\beta_k\theta^k.
\]
Solving gives \eqref{eq:appendix-47}, and direct differentiation gives
\[
\partial_\theta T_1
=-d\theta^{d-1}-\sum_{k=1}^{d-1}k\beta_k\theta^{k-1}.
\]
Since $|\theta|\leq1$ on $E_1$ and $|\beta_k|\leq R$,
\[
\begin{aligned}
|\partial_\theta T_1|
&\leq d|\theta|^{d-1}
 +\sum_{k=1}^{d-1}k|\beta_k||\theta|^{k-1}\\
&\leq d+R\sum_{k=1}^{d-1}k
=d+\frac{Rd(d-1)}2.
\end{aligned}
\]
At $\theta=0$, the $k=1$ term has the ordinary factor
$\theta^0=1$ and all positive powers vanish; this chart contains no negative
power.  At $d=1$, the equation is $\theta+\alpha_0=0$, so the beta tuple and
sum are empty, $T_1=-\theta$, and its speed is one, equal to
$1+R\cdot1\cdot0/2$.
\end{proof}

\begin{lemma}[Highest-lower-degree monic chart]
\label{lem:monic-high-chart}
Under Assumption~\ref{assump:parameter-regime},
Proposition~\ref{prop:monic-presentation}, and
Lemma~\ref{lem:monic-pivot-partition}, if $d\geq2$,
$\theta\in E_d$, and
$\beta=(\beta_0,\ldots,\beta_{d-2})\in[-R,R]^{d-1}$, then
\begin{equation}
T_d(\theta,\beta)
=-\theta-\sum_{k=0}^{d-2}\beta_k\theta^{k-d+1},
\label{eq:appendix-49}
\end{equation}
and
\begin{equation}
|\partial_\theta T_d(\theta,\beta)|
\leq1+R\sum_{k=0}^{d-2}(d-1-k)
=1+\frac{Rd(d-1)}2
\leq d+\frac{Rd(d-1)}2.
\label{eq:appendix-50}
\end{equation}
All negative powers are confined to $E_d\subset\{|\theta|>1\}$.
\end{lemma}

\begin{proof}
Index nonpivot coordinates by exponent.  The active pivot is
$F_d(\theta)=\theta^{d-1}\neq0$, so the root equation
\[
0=\theta^d+\sum_{k=0}^{d-2}\beta_k\theta^k
 +\alpha_{d-1}\theta^{d-1}
\]
can be divided by $\theta^{d-1}$, giving \eqref{eq:appendix-49}.  Differentiating on the
strict high cell gives
\[
\partial_\theta T_d
=-1+\sum_{k=0}^{d-2}(d-1-k)\beta_k\theta^{k-d}.
\]
Here $k-d<0$ and $|\theta|>1$, hence $|\theta^{k-d}|\leq1$.  Thus
\[
\begin{aligned}
|\partial_\theta T_d|
&\leq1+\sum_{k=0}^{d-2}(d-1-k)
 |\beta_k||\theta|^{k-d}\\
&\leq1+R\sum_{k=0}^{d-2}(d-1-k)\\
&=1+\frac{Rd(d-1)}2.
\end{aligned}
\]
Since $d\geq2$, $1\leq d$, proving \eqref{eq:appendix-50}.

At $d=2$, the two formulas reduce visibly to
\[
T_1=-\theta^2-\beta_1\theta,
\qquad
|T_1'|\leq2+R
\quad(|\theta|\leq1),
\]
and
\[
T_2=-\theta-\beta_0\theta^{-1},
\qquad
T_2'=-1+\beta_0\theta^{-2},
\qquad
|T_2'|\leq1+R\leq2+R
\quad(|\theta|>1).
\]
The negative powers are never evaluated at $0$ or $|\theta|=1$; those
points belong to the constant-pivot cell.
\end{proof}

\subsection{Lossless affine-monic specialization}
\label{app:monic-specialization}

\begin{proposition}[Lossless monic transfer through the affine bound]
\label{prop:monic-affine-transfer}
Under Assumptions~\ref{assump:parameter-regime} and
\ref{assump:cube-density-laws}, Proposition~\ref{prop:affine-bound},
Proposition~\ref{prop:monic-presentation}, and
Lemma~\ref{lem:monic-pivot-partition}, let $d\geq1$,
$\mu\in\mathcal D_{d,R,\kappa}$, and let $J\subset\mathbb R$ be a bounded
interval.  For the monic presentation and pivot cells in
Proposition~\ref{prop:monic-presentation} and
Lemma~\ref{lem:monic-pivot-partition},
\begin{equation}
\begin{aligned}
&\Pr_{\alpha\sim\mu}
[\exists\theta\in J:p_\alpha(\theta)=0]\\
&\quad\leq\kappa\sum_{j=1}^d
\int_{E_j}\int_{[-R,R]^{d-1}}
|\partial_\theta T_j(\theta,\beta)|\,d\beta\,d\theta.
\end{aligned}
\label{eq:appendix-51}
\end{equation}
The random vector has exactly $d$ coordinates, all lower coefficients, and
the leading monic coefficient is deterministic.
\end{proposition}

\begin{proof}
Fix $d$, then an arbitrary $\mu\in\mathcal D_{d,R,\kappa}$, and then an
arbitrary bounded $J$.  Proposition~\ref{prop:monic-presentation} supplies a
nondegenerate deterministic $\Theta\supseteq J$, sets $N=d$, and proves the
pointwise identity
\begin{equation}
F_0(\theta)+\langle\alpha,F(\theta)\rangle
=\theta^d+\sum_{k=0}^{d-1}\alpha_k\theta^k
=p_\alpha(\theta).
\label{eq:appendix-52}
\end{equation}
Thus the affine event equals, rather than contains, the polynomial-root
event.  The same proposition supplies deterministic $C^1$ data and keeps
$\alpha=(\alpha_0,\ldots,\alpha_{d-1})\in\mathbb R^d$ without a leading
coordinate.

Lemma~\ref{lem:monic-pivot-partition} supplies the measurable partition of
$J$ and proves that every selected pivot is nonzero.  All interval,
dimension, object, regularity, and pivot hypotheses of
Proposition~\ref{prop:affine-bound} therefore hold with $I=J$ and $N=d$.
Assumption~\ref{assump:cube-density-laws} supplies its one full
$d$-dimensional density cap, with arbitrary correlation.  Applying that
proposition gives \eqref{eq:appendix-51}.  No chart union bound, coordinate conditioning,
law factorization, or root count is added; cube faces and endpoint
conventions remain in the affine interface.
\end{proof}

\begin{lemma}[Coefficient-one two-cell ledger]
\label{lem:monic-two-cell-ledger}
Under Assumption~\ref{assump:parameter-regime} and
Lemmas~\ref{lem:monic-pivot-partition},
\ref{lem:monic-low-chart}, and
\ref{lem:monic-high-chart}, if $d\geq2$ and $J$ is bounded, then
\begin{equation}
\sum_{j=1}^d\int_{E_j}\int_{[-R,R]^{d-1}}
|\partial_\theta T_j(\theta,\beta)|\,d\beta\,d\theta
\leq(2R)^{d-1}
\left(d+\frac{Rd(d-1)}2\right)|J|.
\label{eq:appendix-53}
\end{equation}
The sharper high-chart speed is retained until its coefficient-one
domination; no chart-count factor occurs.
\end{lemma}

\begin{proof}
Define within this proof
\[
V_d:=d+\frac{Rd(d-1)}2,
\qquad
W_d:=1+\frac{Rd(d-1)}2.
\]
The closed beta cube has the exact product volume
\begin{equation}
\operatorname{Leb}^{d-1}([-R,R]^{d-1})
=\prod_{\ell=1}^{d-1}(2R)=(2R)^{d-1}.
\label{eq:appendix-54}
\end{equation}
Cube faces are included and do not alter either this identity or the
pointwise velocity estimates.

All cells other than $E_1$ and $E_d$ are empty.  The two chart lemmas give
\begin{equation}
\int_{E_1}\int_{[-R,R]^{d-1}}
|\partial_\theta T_1|\,d\beta\,d\theta
\leq(2R)^{d-1}V_d|E_1|,
\label{eq:appendix-55}
\end{equation}
\begin{equation}
\int_{E_d}\int_{[-R,R]^{d-1}}
|\partial_\theta T_d|\,d\beta\,d\theta
\leq(2R)^{d-1}W_d|E_d|.
\label{eq:appendix-56}
\end{equation}
Neither cell length has been replaced by $|J|$.  Since $d\geq2$,
\begin{equation}
V_d-W_d=d-1\geq0,
\qquad W_d\leq V_d.
\label{eq:appendix-57}
\end{equation}
Therefore
\begin{equation}
\begin{aligned}
&\sum_{j=1}^d\int_{E_j}\int_{[-R,R]^{d-1}}
|\partial_\theta T_j|\,d\beta\,d\theta\\
&\quad\leq(2R)^{d-1}
 (V_d|E_1|+W_d|E_d|)\\
&\quad\leq(2R)^{d-1}V_d(|E_1|+|E_d|).
\end{aligned}
\label{eq:appendix-58}
\end{equation}
The non-strict low condition assigns both transition points to $E_1$ and the
strict complementary condition assigns every other point to $E_d$.  Hence,
for every endpoint convention,
\begin{equation}
|E_1|+|E_d|=|J|.
\label{eq:appendix-59}
\end{equation}
Substituting \eqref{eq:appendix-59} in \eqref{eq:appendix-58} proves \eqref{eq:appendix-53}.  If either cell is empty, its
integral and length vanish.  The calculation is independent of the location
of $J$.
\end{proof}

\begin{proposition}[Zero-dimensional linear branch]
\label{prop:monic-linear-branch}
Under Assumptions~\ref{assump:parameter-regime} and
\ref{assump:cube-density-laws}, Proposition~\ref{prop:monic-affine-transfer},
Proposition~\ref{prop:monic-presentation}, and
Lemmas~\ref{lem:monic-pivot-partition} and
\ref{lem:monic-low-chart}, if $d=1$,
$\mu\in\mathcal D_{1,R,\kappa}$, and $J$ is bounded, then
\begin{equation}
\Pr_{\alpha_0\sim\mu}
[\exists\theta\in J:\theta+\alpha_0=0]
\leq\kappa|J|
=\kappa(2R)^0
 \left(1+\frac{R\cdot1\cdot0}{2}\right)|J|.
\label{eq:appendix-60}
\end{equation}
\end{proposition}

\begin{proof}
At $d=1$, Proposition~\ref{prop:monic-presentation} gives $N=1$,
$\alpha=(\alpha_0)$, and $p_\alpha(\theta)=\theta+\alpha_0$, with a
deterministic coefficient of $\theta$.  The prescribed partition is
$E_1=J$, and there are no nonpivot coordinates:
\[
[-R,R]^0=\{()\},
\qquad
\operatorname{Leb}^0([-R,R]^0)=1=(2R)^0.
\]
Lemma~\ref{lem:monic-low-chart} gives
$T_1(\theta)=-\theta$ and $|T_1'(\theta)|=1$.  Therefore
\begin{equation}
\int_{E_1}\int_{[-R,R]^0}|T_1'(\theta)|
\,d\operatorname{Leb}^0(\beta)\,d\theta
=\int_J1\,d\theta=|J|.
\label{eq:appendix-61}
\end{equation}
Proposition~\ref{prop:monic-affine-transfer} gives the first inequality in
\eqref{eq:appendix-60}, and $1+R\cdot1\cdot0/2=1$ gives the equality.  There is no high
chart, second cell, or positive-dimensional cube factor.  Empty and
singleton intervals are included by the affine interface.
\end{proof}

\begin{proposition}[Exact affine-monic probability bound]
\label{prop:monic-probability}
Under Assumptions~\ref{assump:parameter-regime} and
\ref{assump:cube-density-laws}, Propositions
~\ref{prop:monic-affine-transfer} and
\ref{prop:monic-linear-branch}, and
Lemma~\ref{lem:monic-two-cell-ledger}, every $d\geq1$, every
$\mu\in\mathcal D_{d,R,\kappa}$, and every bounded interval
$J\subset\mathbb R$ satisfy
\begin{equation}
\Pr_{\alpha\sim\mu}
[\exists\theta\in J:p_\alpha(\theta)=0]
\leq\kappa(2R)^{d-1}
\left(d+\frac{Rd(d-1)}2\right)|J|.
\label{eq:appendix-62}
\end{equation}
The leading coefficient remains deterministic, and arbitrary correlation
among the lower coefficients is allowed.
\end{proposition}

\begin{proof}
Fix $d$, then an arbitrary admissible law, and then a bounded interval with
its literal endpoint convention.  Proposition
~\ref{prop:monic-affine-transfer} applies the general affine inequality to
the identical monic event, the exact $d$-dimensional lower-coefficient
vector, and the prescribed cells.

If $d\geq2$, Lemma~\ref{lem:monic-two-cell-ledger} bounds its complete chart
sum by the right side of \eqref{eq:appendix-53}.  Its calculation first uses the sharper
high-chart speed, then \eqref{eq:appendix-57}, and only then the disjoint length ledger
\eqref{eq:appendix-59}, so it introduces no factor for the number of charts.  Multiplication
by the literal affine coefficient $\kappa$ proves \eqref{eq:appendix-62}.

If $d=1$, Proposition~\ref{prop:monic-linear-branch} gives the same formula
from zero-dimensional volume one and speed one.  These cases exhaust all
$d\geq1$.  The result uses ordinary probability, the original joint density,
and the deterministic leading coefficient; it contains no independent root
theorem, probability-mode conversion, auxiliary threshold, conservative
remainder, or location-dependent constant.
\end{proof}

\subsection{Counter-example presentation and scale certificate}
\label{app:counter-example}

\begin{proposition}[Counter-example presentation and projective speed]
\label{prop:counter-presentation}
Under Assumptions~\ref{assump:parameter-regime},
\ref{assump:balcan-common-chain},
\ref{assump:anchored-unit-range}, and
\ref{assump:cube-density-laws}, suppose
\[
0<\delta\leq1,
\qquad
\Theta=[-1,1],
\qquad
Q_1(x)=1,
\qquad
Q_2(x)=x/\delta,
\]
with $c=0$, $h=1$, $q=0$, and the uniform law on $[-1,1]^2$.  Then
\[
F(\theta)=\left(1,\frac{\theta}{\delta}\right),
\qquad
\langle\alpha,F(\theta)\rangle
=\alpha_1F_1(\theta)+\alpha_2F_2(\theta),
\]
and
\begin{equation}
q=0,
\quad M=0,
\quad B_P=0,
\quad \Delta=1,
\quad N=2,
\quad R=1,
\quad \kappa=\frac14,
\quad A=1,
\quad B_Q=\frac1\delta,
\quad h=1.
\label{eq:appendix-63}
\end{equation}
Moreover,
\begin{equation}
\gamma_F(\theta)
=\frac{(\delta,\theta)}{\sqrt{\delta^2+\theta^2}},
\qquad
\gamma_F'(\theta)
=\frac{(-\delta\theta,\delta^2)}
       {(\delta^2+\theta^2)^{3/2}},
\label{eq:appendix-64}
\end{equation}
and
\begin{equation}
\|\gamma_F'(\theta)\|_2
=\frac{\delta}{\delta^2+\theta^2},
\qquad
\Gamma_{\mathrm{proj}}(F)=\frac1\delta.
\label{eq:appendix-65}
\end{equation}
\end{proposition}

\begin{proof}
Here $x(\theta)=\theta$.  Since $q=0$, the conventions give
$M=B_P=0$.  The two coefficient norms are
\[
\|\operatorname{coeff}(Q_1)\|_1=1,
\qquad
\|\operatorname{coeff}(Q_2)\|_1=\frac1\delta,
\]
and $0<\delta\leq1$ implies $B_Q=1/\delta$.  The output degrees are zero
and one, so $\Delta=1$, and $N=2$.  The literal first feature leaves the
original coefficient ordering unchanged.

The density
\[
f_\square(a_1,a_2)
=\frac14\mathbf1_{[-1,1]^2}(a_1,a_2)
\]
is supported on the square, has essential supremum $1/4$, and integrates to
one.  Thus $\mu_\square\in\mathcal D_{2,1,1/4}$ and
\[
A=(2R)^N\kappa=2^2\cdot\frac14=1,
\]
proving \eqref{eq:appendix-63}.

Multiplying numerator and denominator of $F/\|F\|_2$ by $\delta>0$ gives
the first identity in \eqref{eq:appendix-64}, and direct differentiation gives the second.
Consequently,
\[
\begin{aligned}
\|\gamma_F'(\theta)\|_2
&=\frac{\sqrt{\delta^2\theta^2+\delta^4}}
        {(\delta^2+\theta^2)^{3/2}}\\
&=\frac{\delta\sqrt{\delta^2+\theta^2}}
        {(\delta^2+\theta^2)^{3/2}}
=\frac{\delta}{\delta^2+\theta^2}.
\end{aligned}
\]
This continuous speed is maximized at $\theta=0$, where it is $1/\delta$;
its maximum and essential supremum coincide.  At $\theta=\pm1$ its value is
$\delta/(\delta^2+1)$, so no endpoint or finiteness condition is hidden.
\end{proof}

\begin{lemma}[Exact closed coefficient wedges]
\label{lem:counter-wedges}
Under the assumptions of Proposition~\ref{prop:counter-presentation}, let
$0<\epsilon\leq\delta\leq1$ and set
\[
t:=\frac\epsilon\delta\in(0,1].
\]
Define
\begin{equation}
W_+(t):=\{(a_1,a_2)\in[-1,1]^2:
0\leq a_2\leq1,\ -ta_2\leq a_1\leq0\},
\label{eq:appendix-66}
\end{equation}
\begin{equation}
W_-(t):=\{(a_1,a_2)\in[-1,1]^2:
-1\leq a_2\leq0,\ 0\leq a_1\leq-ta_2\}.
\label{eq:appendix-67}
\end{equation}
Then
\begin{equation}
\{a\in[-1,1]^2:\exists\theta\in[0,\epsilon],
 \ a_1+a_2\theta/\delta=0\}
=W_+(t)\cup W_-(t),
\label{eq:appendix-68}
\end{equation}
$W_+(t)\cap W_-(t)=\{(0,0)\}$, and
\begin{equation}
\Pr_{\alpha\sim\mu_\square}
[\exists\theta\in[0,\epsilon]:
 \alpha_1+\alpha_2\theta/\delta=0]
=\frac{\epsilon}{4\delta}.
\label{eq:appendix-69}
\end{equation}
\end{lemma}

\begin{proof}
For $\theta\in[0,\epsilon]$, put $s=\theta/\delta$.  Then
$s\in[0,t]$.  If $a_2\neq0$, the unique candidate is
\[
s_*=-\frac{a_1}{a_2},
\]
and it lies in the interval precisely when
\begin{equation}
0\leq-\frac{a_1}{a_2}\leq t.
\label{eq:appendix-70}
\end{equation}
For $a_2>0$, this becomes $-ta_2\leq a_1\leq0$; for $a_2<0$, it becomes
$0\leq a_1\leq-ta_2$.  These are the two wedges.

The coefficient axes make \eqref{eq:appendix-68} exact, not only almost everywhere.  If
$a_1=0$ and $a_2\neq0$, the root is the included left endpoint
$\theta=0$, and the appropriate closed wedge contains it.  If $a_2=0$ and
$a_1\neq0$, no root exists and both wedges exclude it.  At the origin the
combination vanishes for every $\theta$, and both wedges contain it.  For
nonzero opposite-sign coefficients, $|a_1|=t|a_2|$ is precisely the right
endpoint $\theta=\epsilon$; strict inequality gives an interior root.
Same-sign nonzero coefficients have no root for $s\geq0$.

The complete triangle boundaries are
\begin{equation}
\begin{aligned}
\partial W_+(t)
={}&\{(0,u):0\leq u\leq1\}
\cup\{(-tu,u):0\leq u\leq1\}\\
&\cup\{(v,1):-t\leq v\leq0\},
\end{aligned}
\label{eq:appendix-71}
\end{equation}
\begin{equation}
\begin{aligned}
\partial W_-(t)
={}&\{(0,u):-1\leq u\leq0\}
\cup\{(-tu,u):-1\leq u\leq0\}\\
&\cup\{(v,-1):0\leq v\leq t\}.
\end{aligned}
\label{eq:appendix-72}
\end{equation}
They include the coefficient-axis edges, roots at $\epsilon$, square-support
edges, and all vertices.  The triangles can overlap only at $a_2=0$, where
their inequalities force $a_1=0$.  Each boundary is a finite union of line
segments and the overlap is one point, so all are planar-null while remaining
part of the exact event.

Direct integration gives
\[
\operatorname{Leb}^2(W_+(t))=\int_0^1tu\,du=\frac t2,
\qquad
\operatorname{Leb}^2(W_-(t))=\int_{-1}^0t|u|\,du=\frac t2.
\]
The null overlap gives
\begin{equation}
\operatorname{Leb}^2(W_+(t)\cup W_-(t))=t.
\label{eq:appendix-73}
\end{equation}
Multiplication by the constant density $1/4$ proves \eqref{eq:appendix-69}.
When $\epsilon=\delta$, $t=1$, each wedge has area $1/2$, the union has
area one, and the probability is $1/4$.  When $\delta=1$, the probability
is $\epsilon/4$ for every $0<\epsilon\leq1$.  Thus all parameter boundaries
and both interval endpoints are included without a limit.
\end{proof}

\begin{proposition}[Counter-example lower certificate and upper scales]
\label{prop:counter-scale}
Under Assumptions~\ref{assump:parameter-regime},
\ref{assump:balcan-common-chain},
\ref{assump:anchored-unit-range}, and
\ref{assump:cube-density-laws}, Proposition~\ref{prop:central-rate},
Proposition~\ref{prop:counter-presentation}, and
Lemma~\ref{lem:counter-wedges}, the Counter-example specialization with
$\mathcal D=\mathcal D_{2,1,1/4}$ satisfies
\begin{equation}
\frac1{4\delta}
\leq C^{\mathrm{Pf}}_{\mathcal D}(F;[-1,1])
\leq\frac1\delta
\leq\frac{\sqrt2}{\delta}.
\label{eq:appendix-74}
\end{equation}
The middle value is the central projective coefficient and the final value is
the raw-presentation coefficient.  No equality or optimality conclusion for
the capacity is asserted.
\end{proposition}

\begin{proof}
Fix any $0<\epsilon\leq\delta$.  Proposition
~\ref{prop:counter-presentation} gives
$\mu_\square\in\mathcal D_{2,1,1/4}$, and
$[0,\epsilon]\subseteq[-1,1]$ has length $\epsilon$.  By
Lemma~\ref{lem:counter-wedges},
\begin{equation}
\frac{
\Pr_{\alpha\sim\mu_\square}
[\exists\theta\in[0,\epsilon]:
 \langle\alpha,F(\theta)\rangle=0]}
{|[0,\epsilon]|}
=\frac{\epsilon/(4\delta)}\epsilon
=\frac1{4\delta}.
\label{eq:appendix-75}
\end{equation}
The capacity supremum contains this law--interval pair, so
\begin{equation}
C^{\mathrm{Pf}}_{\mathcal D}(F;[-1,1])
\geq\frac1{4\delta}.
\label{eq:appendix-76}
\end{equation}

For the upper comparison, Proposition~\ref{prop:central-rate} and the tuple
and speed in Proposition~\ref{prop:counter-presentation} give
\begin{equation}
A\sqrt{\frac N2}\,\Gamma_{\mathrm{proj}}(F)
=1\cdot\sqrt{\frac22}\cdot\frac1\delta
=\frac1\delta,
\label{eq:appendix-77}
\end{equation}
\begin{equation}
\frac{AN\Delta B_Q(1+qB_P)}{\sqrt2\,h}
=\frac{1\cdot2\cdot1\cdot(1/\delta)\cdot1}
       {\sqrt2\cdot1}
=\frac{\sqrt2}{\delta}.
\label{eq:appendix-78}
\end{equation}
Combining \eqref{eq:appendix-76}--\eqref{eq:appendix-78} proves \eqref{eq:appendix-74}.  At $\delta=1$, the three scales
are $1/4$, $1$, and $\sqrt2$; at $\epsilon=\delta$, the witness probability
is $1/4$ and the quotient remains $1/(4\delta)$.  Neither boundary changes
the law, coefficient convention, deterministic feature, ordinary-probability
mode, or Euclidean norm.
\end{proof}

\subsection{Joint central, affine, monic, and scale conclusions}
\label{app:full-conjunction}

\begin{proposition}[Joint anchored sweep conclusions]
\label{prop:full-conjunction}
Under Assumptions~\ref{assump:parameter-regime},
\ref{assump:balcan-common-chain}, and
\ref{assump:anchored-unit-range}, with
Assumption~\ref{assump:cube-density-laws} in every probabilistic clause and
Assumption~\ref{assump:affine-chart-data} in the general affine clause, the
following conclusions hold simultaneously.

First, with $D_*=\Delta B_Q(1+qB_P)$,
\[
|G_i'(x)|\leq D_*,
\qquad
\|G'(x)\|_2\leq\sqrt N\,D_*,
\]
$G_1=F_1=1$, and
\begin{equation}
\gamma_G'
=\frac{(I_N-\gamma_G\gamma_G^{\mathsf T})G'}{\|G\|_2},
\qquad
\gamma_F'(\theta)
=\frac1h\gamma_G'\!\left(\frac{\theta-c}{h}\right),
\qquad
\Gamma_{\mathrm{proj}}(F)
\leq\frac{\sqrt N\,D_*}{h}.
\label{eq:appendix-79}
\end{equation}
This includes one-sided endpoints, $q=0$, $N=1$, $\Delta=0$, and stationary
normalized curves, with degree-zero dependence on $M$ for fixed $B_P$.

Second, after fixing the deterministic presentation, every
$\mu\in\mathcal D_{N,R,\kappa}$ and then every positive-length
$I\subseteq\Theta$ satisfy
\begin{equation}
\Pr_{\alpha\sim\mu}
[\exists\theta\in I:\langle\alpha,F(\theta)\rangle=0]
\leq A\sqrt{\frac N2}\,\Gamma_{\mathrm{proj}}(F)|I|
\leq\frac{AN\Delta B_Q(1+qB_P)}{\sqrt2 h}|I|,
\label{eq:appendix-80}
\end{equation}
and
\begin{equation}
C^{\mathrm{Pf}}_{\mathcal D}(F;\Theta)
\leq A\sqrt{\frac N2}\,\Gamma_{\mathrm{proj}}(F)
\leq\frac{AN\Delta B_Q(1+qB_P)}{\sqrt2 h}.
\label{eq:appendix-81}
\end{equation}
The capacity retains the interval-then-law supremum order.  Before division,
empty and singleton intervals have zero right side.  All endpoint,
multiplicity, stationary, and identically-zero cases are included.

Third, for every affine invocation and every
$\mu\in\mathcal D_{N,R,\kappa}$,
\begin{equation}
\begin{aligned}
&\Pr_{\alpha\sim\mu}
[\exists\theta\in I:F_0(\theta)+\langle\alpha,F(\theta)\rangle=0]\\
&\quad\leq\kappa\sum_{j=1}^N
\int_{E_j}\int_{[-R,R]^{N-1}}
|\partial_\theta T_j(\theta,\beta)|\,d\beta\,d\theta,
\end{aligned}
\label{eq:appendix-82}
\end{equation}
in $[0,+\infty]$.  The original coefficient dimension, arbitrary
correlation, near-zero pivots, all root multiplicities, $N=1$, and divergent
integrals remain admissible.

Fourth, for every $d\geq1$ and bounded $J\subset\mathbb R$, there is a
nondegenerate $\Theta=[c-h,c+h]\supseteq J$ with
\[
Q_0(x)=(c+hx)^d,
\quad Q_{k+1}(x)=(c+hx)^k,
\quad F_0(\theta)=\theta^d,
\quad F_{k+1}(\theta)=\theta^k,
\]
\begin{equation}
p_\alpha(\theta)=\theta^d+\sum_{k=0}^{d-1}\alpha_k\theta^k,
\qquad
(q,M,B_P,N,A,\Delta_{\mathrm{aug}})
=(0,0,0,d,(2R)^d\kappa,d).
\label{eq:appendix-83}
\end{equation}
The monic coefficient is deterministic.  For $d\geq2$, the prescribed cells
are those in \eqref{eq:appendix-46}, and
\begin{equation}
T_1=-\theta^d-\sum_{k=1}^{d-1}\beta_k\theta^k,
\qquad
|\partial_\theta T_1|
\leq d+\frac{Rd(d-1)}2
\quad(|\theta|\leq1),
\label{eq:appendix-84}
\end{equation}
\begin{equation}
T_d=-\theta-\sum_{k=0}^{d-2}\beta_k\theta^{k-d+1},
\qquad
|\partial_\theta T_d|
\leq1+\frac{Rd(d-1)}2
\leq d+\frac{Rd(d-1)}2
\quad(|\theta|>1).
\label{eq:appendix-85}
\end{equation}
At $d=1$, $E_1=J$, $T_1=-\theta$, and $|T_1'|=1$.  Every possibly
correlated $\mu\in\mathcal D_{d,R,\kappa}$ satisfies
\begin{equation}
\Pr_{\alpha\sim\mu}
[\exists\theta\in J:p_\alpha(\theta)=0]
\leq\kappa(2R)^{d-1}
\left(d+\frac{Rd(d-1)}2\right)|J|.
\label{eq:appendix-86}
\end{equation}

Fifth, for $0<\delta\leq1$, the uniform-square specialization has
\[
(h,q,M,\Delta,N,R,\kappa,A,B_P,B_Q)
=\left(1,0,0,1,2,1,\frac14,1,0,\frac1\delta\right),
\]
\begin{equation}
\gamma_F(\theta)
=\frac{(\delta,\theta)}{\sqrt{\delta^2+\theta^2}},
\qquad
\|\gamma_F'(\theta)\|_2
=\frac{\delta}{\delta^2+\theta^2},
\qquad
\Gamma_{\mathrm{proj}}(F)=\frac1\delta.
\label{eq:appendix-87}
\end{equation}
For every $0<\epsilon\leq\delta$,
\begin{equation}
\Pr[\exists\theta\in[0,\epsilon]:
 \alpha_1+\alpha_2\theta/\delta=0]
=\frac{\epsilon}{4\delta},
\label{eq:appendix-88}
\end{equation}
and therefore
\begin{equation}
\frac1{4\delta}
\leq C^{\mathrm{Pf}}_{\mathcal D_{2,1,1/4}}(F;[-1,1])
\leq\frac1\delta
\leq\frac{\sqrt2}{\delta}.
\label{eq:appendix-89}
\end{equation}
The last chain does not claim equality or optimality for the capacity.

All displayed constants are literal and have no hidden dependence.  The
probability mode is ordinary probability under a fixed full joint law; the
central horizon mode is per interval followed only where displayed by the
interval supremum and then the law supremum.  Central vector and operator
norms are Euclidean, chart velocities are scalar, and coefficient and
interval measures use the conventions in Section~\ref{sec:prelim}.  No confidence
parameter, independence assumption, asymptotic horizon, or normalization
beyond the declared anchored setting is present.
\end{proposition}

\begin{proof}
Proposition~\ref{prop:coordinate-envelope} gives the coordinatewise part of
\eqref{eq:appendix-79}, including the endpoint and degree-zero branches.
Proposition~\ref{prop:projective-speed} gives anchor-derived nonvanishing,
the normalized-projector identity, the exact coordinate change, and the
remaining projective-speed assertions in \eqref{eq:appendix-79}.

Proposition~\ref{prop:central-rate} gives \eqref{eq:appendix-80}--\eqref{eq:appendix-81} with the same
deterministic feature, ordinary-probability mode, literal constants, and
ordered suprema.  Proposition~\ref{prop:affine-bound} gives \eqref{eq:appendix-82} in the
original coefficient dimension and extended-real mode.

Proposition~\ref{prop:monic-presentation} gives the objects and tuple in
\eqref{eq:appendix-83}.  Lemma~\ref{lem:monic-pivot-partition} gives the prescribed cells,
and Lemmas~\ref{lem:monic-low-chart} and
\ref{lem:monic-high-chart} give \eqref{eq:appendix-84}--\eqref{eq:appendix-85}, including the transition
behavior and the zero-dimensional convention.  Proposition
~\ref{prop:monic-probability} supplies only the probability conclusion
\eqref{eq:appendix-86}, on the same monic object and lower-coefficient law.

Proposition~\ref{prop:counter-presentation} gives the tuple and speed in
\eqref{eq:appendix-87}, Lemma~\ref{lem:counter-wedges} gives \eqref{eq:appendix-88} with all closed
boundaries, and Proposition~\ref{prop:counter-scale} gives \eqref{eq:appendix-89}.  Each
clause is therefore a conclusion of the named result in the identical
object, dimension, law, interval, norm, and probability mode.  Their logical
conjunction proves the proposition without an additional estimate,
assumption, probability conversion, or mode change.
\end{proof}

\subsection{Proof of the Main Theorem}

\begin{proof}[Proof of Theorem~\ref{thm:main}]
The primitive derivative and projective clauses follow from
Propositions~\ref{prop:coordinate-envelope} and
\ref{prop:projective-speed}.  The central probability and capacity clauses,
including the literal algebraic bridge and the interval-then-law supremum
order, are Proposition~\ref{prop:central-rate}.  The general affine clause is
Proposition~\ref{prop:affine-bound}.

For the detailed monic interface, use
Proposition~\ref{prop:monic-presentation},
Lemma~\ref{lem:monic-pivot-partition}, and
Lemmas~\ref{lem:monic-low-chart} and
\ref{lem:monic-high-chart}; the exact monic probability rate is
Proposition~\ref{prop:monic-probability}, whose two dimensional branches are
implemented by Lemma~\ref{lem:monic-two-cell-ledger} and
Proposition~\ref{prop:monic-linear-branch}.

For Counter-example 1, Proposition~\ref{prop:counter-presentation} supplies
the tuple and Euclidean speed, Lemma~\ref{lem:counter-wedges} supplies the
closed-event probability, and Proposition~\ref{prop:counter-scale} supplies
the three distinct scale comparisons.  Proposition
~\ref{prop:full-conjunction} verifies that these named conclusions have the
same assumption bases, objects, dimensions, law and interval scopes, norm
conventions, and ordinary-probability mode as the clauses in
Theorem~\ref{thm:main}.
Their direct conjunction is precisely Theorem~\ref{thm:main}.
\end{proof}
